% ============================================================
%  JUEGOS DIDÁCTICOS — CONJUNTO DE LOS NÚMEROS ENTEROS (Z)
%  Curso: 1er Año — Educación Media General
%  Documento para el docente: fabricación y aplicación en el aula
% ============================================================

\documentclass[12pt, letterpaper]{article}

\usepackage{mathwizards-guia}
\usepackage{longtable}

\mwguia{Juegos Didácticos — Números Enteros}{Matemáticas $\cdot$ 1er Año}

\newcommand{\Z}{\mathbb{Z}}
\newcommand{\abs}[1]{\lvert #1\rvert}
\setlength{\parindent}{0pt}
\setlength{\parskip}{4pt}
\renewcommand{\arraystretch}{1.18}

\begin{document}

\mwportada{Juegos Didácticos — Números Enteros}{Quince juegos para el aula}{Orden $\cdot$ Valor absoluto $\cdot$ Operaciones $\cdot$ Propiedades $\cdot$ Agrupación}

\vspace{4pt}
\begin{cuadroinfo}
  \small
  \textbf{Presentación.} Este documento reúne quince juegos para enseñar el conjunto de los números enteros $\Z$ en 1er año. Cada juego sigue \emph{siempre la misma estructura}: ficha rápida, idea central, resumen en cinco pasos, reglas, turno paso a paso, final y puntuación, ejemplo resuelto, materiales, construcción, \textbf{patrones explícitos de fabricación} y configuración inicial. Todos los mazos, tableros, dados y tarjetas están resueltos y verificados: el docente solo recorta y juega.\\[4pt]
  \textbf{Tres niveles.} Los \textbf{juegos 1 a 5} son de iniciación: una sola idea por turno, materiales escolares básicos y partidas de 10 a 20 minutos. Los \textbf{juegos 6 a 10} son de entrada: una comparación, una operación o, como máximo, dos pasos por turno. Los \textbf{juegos 11 a 15} profundizan el mismo contenido con más reglas, más piezas y auditoría entre compañeros.\\[4pt]
  \textbf{Regla de oro:} el estudiante debe \emph{decir en voz alta la operación y el resultado} antes de mover una pieza. El juego evalúa la justificación, no la velocidad.
\end{cuadroinfo}

\setcounter{tocdepth}{1}
\newpage
\tableofcontents
\newpage

% ============================================================
\section*{Cómo usar este documento}
\addcontentsline{toc}{section}{Cómo usar este documento}
% ============================================================

Cada juego se presenta con los mismos diez bloques, para que el docente pueda preparar cualquiera de los quince sin releer todo:

\begin{enumerate}[leftmargin=*]
  \item \textbf{Ficha rápida:} jugadores, duración, objetivo y subtemas.
  \item \textbf{La idea en una frase:} el «gancho» que se dice a los estudiantes.
  \item \textbf{Cómo se juega en cinco pasos:} el resumen que cabe en el pizarrón.
  \item \textbf{Reglas:} numeradas y sin ambigüedades.
  \item \textbf{Turno paso a paso:} el procedimiento exacto del jugador.
  \item \textbf{Final y puntuación:} cuándo termina y cómo se gana.
  \item \textbf{Ejemplo resuelto:} una jugada completa con los cálculos a la vista.
  \item \textbf{Materiales y construcción:} qué comprar, qué reciclar y cómo armarlo.
  \item \textbf{Patrones explícitos:} todos los mazos, tableros y tarjetas ya resueltos, más la regla para fabricar más.
  \item \textbf{Configuración inicial:} cómo se dispone la mesa antes de la primera jugada.
\end{enumerate}

\section*{Mapa general de los quince juegos}
\addcontentsline{toc}{section}{Mapa general de los quince juegos}

\subsection*{Juegos de iniciación (1 a 5)}

\begin{center}
\footnotesize
\setlength{\tabcolsep}{4pt}
\begin{tabular}{@{}c l p{3.6cm} p{4.6cm} c c@{}}
\toprule
\textbf{\#} & \textbf{Juego} & \textbf{Subtemas de $\Z$} & \textbf{Mecánica central} & \textbf{Jug.} & \textbf{Tiempo}\\
\midrule
1 & \textbf{Semáforo de enteros} & Orden, opuestos, valor absoluto, operaciones & Verdadero o falso con tarjetas de color & 2--4 o clase & 10--15 min\\
2 & \textbf{El termómetro humano} & Ubicación y orden en la recta & Fila humana sobre una cinta numérica & 2--4 & 15--20 min\\
3 & \textbf{El ascensor de dados} & Suma y resta de una operación & Dado con cargas y ficha en una tira numérica & 2--4 & 15--20 min\\
4 & \textbf{Pesca de opuestos} & Números opuestos & Pesca de cartas entre compañeros & 2--4 & 15--20 min\\
5 & \textbf{Tres en raya de enteros} & Suma y resta de una operación & Cuadrícula $3\times 3$ con resultados & 2 o 2 eq. & 15--20 min\\
\bottomrule
\end{tabular}
\end{center}

\subsection*{Juegos de entrada (6 a 10)}

\begin{center}
\footnotesize
\setlength{\tabcolsep}{4pt}
\begin{tabular}{@{}c l p{3.6cm} p{4.6cm} c c@{}}
\toprule
\textbf{\#} & \textbf{Juego} & \textbf{Subtemas de $\Z$} & \textbf{Mecánica central} & \textbf{Jug.} & \textbf{Tiempo}\\
\midrule
6 & \textbf{Guerra de pisos} & Orden y comparación & Guerra de cartas sobre la recta vertical & 2--4 & 15--20 min\\
7 & \textbf{Memoria de opuestos} & Opuestos y valor absoluto & Memoria de parejas & 2--4 & 15--20 min\\
8 & \textbf{La serpiente de los enteros} & Suma y resta de una operación & Tablero de 30 casillas con saltos & 2--4 & 20--25 min\\
9 & \textbf{Bingo de productos} & Producto, cociente exacto y signos & Bingo de $4\times 4$ con dos mazos & 2--4 & 20--25 min\\
10 & \textbf{La escalera de dos peldaños} & Combinadas de dos pasos & Carrera de 8 peldaños & 2--4 & 20--25 min\\
\bottomrule
\end{tabular}
\end{center}

\subsection*{Juegos de profundización (11 a 15)}

\begin{center}
\footnotesize
\setlength{\tabcolsep}{4pt}
\begin{tabular}{@{}c l p{3.6cm} p{4.6cm} c c@{}}
\toprule
\textbf{\#} & \textbf{Juego} & \textbf{Subtemas de $\Z$} & \textbf{Mecánica central} & \textbf{Jug.} & \textbf{Tiempo}\\
\midrule
11 & \textbf{Abismo y Cumbre} & Orden, opuestos, valor absoluto & Recta numérica y captura por simetría & 2--4 & 30--40 min\\
12 & \textbf{Laboratorio de Antimateria} & Adición, sustracción, propiedades & Reactor de cargas y pares nulos & 2--4 & 25--35 min\\
13 & \textbf{Vectores de Invasión} & Producto, cociente, signos, distributiva & Conquista de una matriz $6\times 6$ & 2 o 2 eq. & 30--45 min\\
14 & \textbf{Cortocircuito $\Z$} & Combinadas sin agrupación & Dominó de expresiones con auditoría & 2--4 & 25--35 min\\
15 & \textbf{Código Bóveda} & Combinadas con agrupación & Tres anillos concéntricos & 2--4 & 30--45 min\\
\bottomrule
\end{tabular}
\end{center}

\textbf{Orden pedagógico sugerido:} $1\to2\to3\to\cdots\to15$. Cada juego supone dominado el anterior. Si el grupo ya domina lo básico, se puede empezar directamente en el Juego 11 (profundización); si necesita un calentamiento, conviene empezar en el Juego 1.

\section*{Acuerdos comunes}
\addcontentsline{toc}{section}{Acuerdos comunes}

\begin{enumerate}[leftmargin=*]
  \item \textbf{Sin calculadora.} Cada jugador tiene una hoja de borrador (bitácora, multímetro o bloc de criptografía).
  \item \textbf{Se canta la jugada.} Antes de mover o colocar una pieza se anuncia la operación y el resultado.
  \item \textbf{El error se corrige, no se castiga dos veces.} Se pierde el turno o la pieza según la regla, pero siempre se explica dónde estuvo el error.
  \item \textbf{Los roles rotan.} En equipos: operador, auditor, cronometrador y escriba.
  \item \textbf{El material se cuida.} Cada mazo se guarda en un sobre rotulado para reutilizarlo en otros cursos.
  \item \textbf{Se puede subir o bajar de nivel.} Los juegos 1 a 5 inician, los juegos 6 a 10 consolidan y los juegos 11 a 15 profundizan. Si un estudiante falla en un juego de profundización, conviene volver al juego de entrada equivalente; si el problema es de base, se vuelve a un juego de iniciación.
\end{enumerate}

% ============================================================
\clearpage
\section{Juego 1 — Semáforo de enteros: ¿verdadero o falso?}
% ============================================================

\subsection{Ficha rápida}

\begin{center}
\small
\begin{tabular}{@{}p{3cm} p{12.2cm}@{}}
\toprule
\textbf{Jugadores} & 2 a 4 (o toda la clase organizada en parejas)\\
\textbf{Duración} & 10--15 minutos\\
\textbf{Objetivo} & Decidir si una afirmación sobre enteros es verdadera o falsa y justificarla con una regla.\\
\textbf{Subtemas} & Orden, opuestos, valor absoluto, suma, resta y signo del producto.\\
\textbf{Qué se evalúa} & Que el estudiante elija la tarjeta correcta y explique la regla en una frase.\\
\textbf{Nivel} & Iniciación (calentamiento de todo el bloque).\\
\bottomrule
\end{tabular}
\end{center}

\subsection{La idea en una frase}

\begin{cuadroextra}
  «Cada afirmación es un semáforo: verde si es verdadera, roja si es falsa. No gana quien levanta primero la tarjeta, sino quien explica por qué.»
\end{cuadroextra}

\subsection{Cómo se juega en cinco pasos}

\begin{enumerate}[leftmargin=*]
  \item Cada jugador tiene una tarjeta verde (V) y una roja (F).
  \item Se lee en voz alta una tarjeta del mazo.
  \item Se cuentan cinco segundos y todos levantan a la vez su tarjeta.
  \item Quien acierta explica la regla en una frase; si no la explica, no suma.
  \item Gana quien sume más puntos al terminar el mazo o al cumplirse el tiempo.
\end{enumerate}

\subsection{Reglas}

\begin{enumerate}[leftmargin=*]
  \item \textbf{Una tarjeta por ronda.} Se lee una sola afirmación y no se repite hasta la ronda siguiente.
  \item \textbf{Silencio hasta la señal.} Nadie habla, señala ni levanta su tarjeta antes de que el docente diga «¡Ya!».
  \item \textbf{Una sola tarjeta.} Se debe levantar exactamente una tarjeta. Quien muestre las dos, o cambie después de la señal, no suma en esa ronda.
  \item \textbf{Acierto con explicación.} Vale \textbf{1 punto} si el jugador acierta y dice la regla que lo justifica.
  \item \textbf{Acierto sin explicación.} Vale \textbf{0 puntos}: la respuesta correcta sin justificación no se considera lograda.
  \item \textbf{Error.} Vale \textbf{0 puntos}. El docente aprovecha la ronda para corregir en la recta o con un ejemplo.
  \item \textbf{Varios aciertos.} Si dos o más jugadores aciertan y justifican, todos suman. No hay competencia por la velocidad.
  \item \textbf{Juez.} El docente (o un estudiante rotativo) lee la tarjeta, controla el tiempo y valida la justificación con la clave del Patrón B.
  \item \textbf{Fin del mazo.} Se puede jugar el mazo completo (24 tarjetas) o fijar un tiempo de 10 minutos; lo que ocurra primero.
\end{enumerate}

\subsection{Turno paso a paso}

\begin{enumerate}[leftmargin=*]
  \item \textbf{Escucha} la afirmación completa sin interrumpir.
  \item \textbf{Piensa} la respuesta y la regla que la justifica.
  \item \textbf{Levanta} una sola tarjeta cuando el docente dé la señal.
  \item \textbf{Explica} la regla en una frase, si acertaste.
  \item \textbf{Anota} tu punto y prepárate para la ronda siguiente.
\end{enumerate}

\subsection{Final y puntuación}

\begin{center}
\small
\begin{tabular}{@{}p{7.4cm} c@{}}
\toprule
\textbf{Concepto} & \textbf{Puntos}\\
\midrule
Respuesta correcta con justificación & $1$\\
Respuesta correcta sin justificación & $0$\\
Respuesta incorrecta & $0$\\
\bottomrule
\end{tabular}
\end{center}

Gana quien sume más puntos al terminar las 24 tarjetas (o los 10 minutos). Desempate: una ronda extra con una tarjeta nueva; si persiste, gana quien haya justificado primero en esa ronda.

\subsection{Ejemplo resuelto}

Se lee: «$-7>-3$».

\begin{cuadrosolucion}
\begin{enumerate}[leftmargin=*]
  \item El jugador piensa: en la recta, $-7$ está a la izquierda de $-3$.
  \item Levanta la tarjeta \textbf{roja (falsa)}.
  \item Explica: «Es falsa, porque $-7$ es menor que $-3$: está más a la izquierda».
  \item Gana $1$ punto.
\end{enumerate}
Otras dos rondas:
\begin{enumerate}[leftmargin=*]
  \item «$\abs{-9}=9$»: \textbf{verdadera}; el valor absoluto es la distancia al cero.
  \item «$-3\cdot 4=+12$»: \textbf{falsa}; menos por más da menos, el resultado es $-12$.
\end{enumerate}
\end{cuadrosolucion}

\subsection{Materiales}

\begin{center}
\small
\begin{tabular}{@{}p{5.0cm} c p{7.4cm}@{}}
\toprule
\textbf{Material} & \textbf{Cantidad} & \textbf{Uso}\\
\midrule
Cartulina verde & 1 pliego & Tarjetas «verdadero»\\
Cartulina roja & 1 pliego & Tarjetas «falso»\\
Cartulina blanca & 2 pliegos & 24 tarjetas de afirmaciones\\
Marcadores negro y rojo & 1 de cada uno & Rotular\\
Tijera, regla y lápiz & 1 de cada uno & Trazado y corte\\
Hoja de puntos & 1 por jugador & Llevar el puntaje\\
\bottomrule
\end{tabular}
\end{center}

\subsection{Construcción paso a paso}

\begin{enumerate}[leftmargin=*]
  \item \textbf{Tarjetas de respuesta.} Recorta cuadrados de $8\text{ cm}\times 8\text{ cm}$: uno verde con una «V» grande y uno rojo con una «F».
  \item \textbf{Tarjetas de afirmación.} Recorta 24 rectángulos de $6\text{ cm}\times 9\text{ cm}$ y copia las afirmaciones del Patrón A.
  \item \textbf{Clave.} En el reverso de cada tarjeta escribe la respuesta (V o F) y la justificación corta.
  \item \textbf{Control.} Verifica que cada afirmación tenga una sola respuesta correcta y que la justificación no admita discusión.
\end{enumerate}

\subsection{Patrones explícitos (fabricación)}

\subsubsection{Patrón A — Las 24 afirmaciones}

{\footnotesize
\begin{longtable}{@{}r >{\raggedright\arraybackslash}p{5.4cm} c >{\raggedright\arraybackslash}p{5.4cm}@{}}
\toprule
\textbf{\#} & \textbf{Afirmación} & \textbf{Resp.} & \textbf{Justificación}\\
\midrule
\endhead
\bottomrule
\endfoot
1 & $-7<-3$ & V & $-7$ está a la izquierda de $-3$ en la recta.\\
2 & $+4>-1$ & V & Todo positivo es mayor que cualquier negativo.\\
3 & $-5>0$ & F & Todo negativo es menor que $0$.\\
4 & $\abs{-9}=9$ & V & El valor absoluto es la distancia al cero.\\
5 & $\abs{-9}=-9$ & F & Una distancia nunca es negativa.\\
6 & El opuesto de $+6$ es $-6$. & V & Cambiar el signo da el opuesto.\\
7 & El opuesto de $-2$ es $-2$. & F & El opuesto de $-2$ es $+2$.\\
8 & $3+(-5)=-2$ & V & Signos distintos: se restan y gana el signo del mayor.\\
9 & $3+(-5)=+8$ & F & No se suman los valores absolutos si los signos son distintos.\\
10 & $-4-(-6)=+2$ & V & Restar un negativo equivale a sumar: $-4+6=+2$.\\
11 & $-4-(-6)=-10$ & F & $-4+6=+2$, no $-10$.\\
12 & $-3\cdot 4=+12$ & F & Menos por más da menos: el resultado es $-12$.\\
13 & $(-2)\cdot(-5)=-10$ & F & Menos por menos da más: el resultado es $+10$.\\
14 & $(-2)\cdot(-5)=+10$ & V & Menos por menos da más.\\
15 & $12\div(-3)=-4$ & V & Más entre menos da menos.\\
16 & $12\div(-3)=+4$ & F & Más entre menos da menos: el resultado es $-4$.\\
17 & $-8<-10$ & F & $-8$ está a la derecha de $-10$.\\
18 & $0>-1$ & V & $0$ está a la derecha de $-1$.\\
19 & $\abs{-5}<\abs{+2}$ & F & $\abs{-5}=5$ y $\abs{+2}=2$; $5$ no es menor que $2$.\\
20 & $5-9=-4$ & V & $5-9=-4$.\\
21 & $5-9=+4$ & F & El resultado es negativo: $-4$.\\
22 & $(-1)\cdot(-1)\cdot(-1)=-1$ & V & Dos negativos dan positivo y el tercero lo vuelve negativo.\\
23 & $-15\div(-5)=-3$ & F & Menos entre menos da más: el resultado es $+3$.\\
24 & El opuesto de $0$ es $0$. & V & El cero es su propio opuesto.\\
\end{longtable}
}

\subsubsection{Patrón B — Regla de puntuación}

\begin{center}
\small
\begin{tabular}{@{}p{8.4cm} c@{}}
\toprule
\textbf{Situación en la ronda} & \textbf{Puntos}\\
\midrule
Tarjeta correcta y regla explicada & $1$\\
Tarjeta correcta sin regla & $0$\\
Tarjeta incorrecta & $0$\\
Muestra las dos tarjetas o cambia después de la señal & $0$\\
\bottomrule
\end{tabular}
\end{center}

\subsubsection{Patrón C — Cómo fabricar nuevas afirmaciones}

\begin{enumerate}[leftmargin=*]
  \item Elige \textbf{una sola idea} por afirmación: orden, opuesto, valor absoluto, suma, resta o signo.
  \item Escribe la afirmación con números pequeños (entre $-15$ y $+15$) para que se resuelva mentalmente.
  \item Comprueba la respuesta con la recta o con la regla de signos antes de escribir el reverso.
  \item Evita las afirmaciones con doble negación o ambigüedad: deben tener una única respuesta defendible.
  \item Equilibra el mazo: al menos 8 verdaderas y 8 falsas, para que no se pueda adivinar por patrón.
\end{enumerate}

\subsection{Configuración inicial}

\begin{enumerate}[leftmargin=*]
  \item Reparte un par de tarjetas (verde y roja) a cada jugador.
  \item Baraja las 24 tarjetas de afirmación y déjalas boca abajo al centro.
  \item Entrega una hoja de puntos y un lápiz a cada jugador.
  \item Acuerda si se juega el mazo completo o 10 minutos.
  \item Designa al juez y al lector de tarjetas (pueden ser el docente y un estudiante).
\end{enumerate}

% ============================================================
\clearpage
\section{Juego 2 — El termómetro humano: la fila de los enteros}
% ============================================================

\subsection{Ficha rápida}

\begin{center}
\small
\begin{tabular}{@{}p{3cm} p{12.2cm}@{}}
\toprule
\textbf{Jugadores} & 2 a 4 (funciona mejor con 4 a 8 en la fila)\\
\textbf{Duración} & 15--20 minutos\\
\textbf{Objetivo} & Ubicar enteros en la recta numérica y ordenarlos usando el cuerpo como recta.\\
\textbf{Subtemas} & Ubicación, orden y comparación de enteros.\\
\textbf{Qué se evalúa} & Que diga el valor, se ubique en la marca correcta y nombre un número mayor y uno menor.\\
\textbf{Nivel} & Iniciación.\\
\bottomrule
\end{tabular}
\end{center}

\subsection{La idea en una frase}

\begin{cuadroextra}
  «La cinta del piso es la recta numérica: el $0$ va al centro, los negativos a la izquierda y los positivos a la derecha. Cada tarjeta es un pasajero que debe subir al bus en el orden correcto.»
\end{cuadroextra}

\subsection{Cómo se juega en cinco pasos}

\begin{enumerate}[leftmargin=*]
  \item Se pega en el piso una cinta con marcas de $-10$ a $+10$ y el $0$ al centro.
  \item Cada jugador recibe tres tarjetas boca abajo.
  \item En tu turno volteas una tarjeta, dices el valor y caminas a su marca.
  \item Si la marca está ocupada, te colocas al lado que corresponde.
  \item Al final de cada ronda se revisa la fila; gana quien sume más puntos en tres rondas.
\end{enumerate}

\subsection{Reglas}

\begin{enumerate}[leftmargin=*]
  \item \textbf{La cinta es la recta.} El $0$ está al centro; los negativos van a la izquierda y los positivos a la derecha. La cinta no se mueve durante la partida.
  \item \textbf{Tres tarjetas por jugador.} Cada jugador recibe tres tarjetas y juega una por ronda.
  \item \textbf{Anuncio obligatorio.} Antes de caminar, el jugador dice en voz alta el valor de su tarjeta.
  \item \textbf{Ubicación exacta.} Se debe pisar la marca del número. Si está ocupada, se coloca a la izquierda si su número es menor, o a la derecha si es mayor.
  \item \textbf{Números iguales.} Si dos jugadores tienen el mismo número, el segundo se coloca junto al primero y ambos dicen «iguales». No se permite ocupar la misma marca.
  \item \textbf{Revisión de ronda.} Cuando todos jugaron una tarjeta, el grupo revisa la fila de izquierda a derecha y comprueba que esté ordenada de menor a mayor.
  \item \textbf{Puntos.} Ubicación correcta: $1$ punto. Además, si nombra un número mayor y uno menor que el suyo: $1$ punto extra. Si se equivoca de marca: $0$ puntos y vuelve a su puesto.
  \item \textbf{Apoyo permitido.} Si un jugador duda, puede mirar la cinta para ubicarse, pero entonces no gana el punto extra por nombrar mayor y menor.
  \item \textbf{Tres rondas.} Al terminar las tres rondas gana quien tenga más puntos. Desempate: quien haya acertado la última ubicación.
\end{enumerate}

\subsection{Turno paso a paso}

\begin{enumerate}[leftmargin=*]
  \item \textbf{Voltea} tu tarjeta y mira el número.
  \item \textbf{Di} el valor en voz alta.
  \item \textbf{Camina} a la marca correspondiente.
  \item \textbf{Acomódate} a la izquierda o a la derecha si la marca está ocupada.
  \item \textbf{Nombra} un número mayor y uno menor que el tuyo.
  \item \textbf{Espera} la revisión de la ronda y anota tus puntos.
\end{enumerate}

\subsection{Final y puntuación}

\begin{center}
\small
\begin{tabular}{@{}p{8.8cm} c@{}}
\toprule
\textbf{Logro} & \textbf{Puntos}\\
\midrule
Ubicación correcta en la marca & $1$\\
Nombrar un mayor y un menor & $1$ extra\\
Ubicación incorrecta & $0$\\
\bottomrule
\end{tabular}
\end{center}

Gana quien sume más puntos tras tres rondas. Si hay empate, gana quien haya acertado la última ubicación; si persiste, se juega una ronda extra con una tarjeta nueva.

\subsection{Ejemplo resuelto}

Al jugador le sale la tarjeta $-4$.

\begin{cuadrosolucion}
\begin{enumerate}[leftmargin=*]
  \item Dice: «Soy $-4$».
  \item Camina cuatro marcas a la izquierda del $0$.
  \item Nombra: «$-7$ es menor que $-4$ y $+2$ es mayor que $-4$».
  \item Gana $2$ puntos: $1$ por la ubicación y $1$ por los dos números.
\end{enumerate}
Otro turno: sale $+6$.
\begin{enumerate}[leftmargin=*]
  \item Se ubica seis marcas a la derecha del $0$.
  \item Nombra: «$-1$ es menor y $+9$ es mayor».
  \item Gana $2$ puntos.
\end{enumerate}
\end{cuadrosolucion}

\subsection{Materiales}

\begin{center}
\small
\begin{tabular}{@{}p{5.0cm} c p{7.4cm}@{}}
\toprule
\textbf{Material} & \textbf{Cantidad} & \textbf{Uso}\\
\midrule
Cinta adhesiva o masking & 1 rollo & Recta numérica en el piso\\
Cartulina & 2 pliegos & 20 tarjetas\\
Marcadores negro y rojo & 1 de cada uno & Positivos y negativos\\
Tijera, regla y lápiz & 1 de cada uno & Trazado y corte\\
Hoja de puntos & 1 por jugador & Llevar el puntaje\\
\bottomrule
\end{tabular}
\end{center}

\subsection{Construcción paso a paso}

\begin{enumerate}[leftmargin=*]
  \item \textbf{Recta.} Extiende $4$ o $5$ metros de cinta en el piso y marca $21$ puntos separados por unos $20$ cm.
  \item \textbf{Numeración.} Escribe $0$ en el centro, $-1$ a $-10$ hacia la izquierda y $+1$ a $+10$ hacia la derecha.
  \item \textbf{Tarjetas.} Recorta 20 rectángulos de $6\text{ cm}\times 9\text{ cm}$ y copia un valor distinto en cada uno (Patrón A).
  \item \textbf{Colores.} Escribe los negativos en rojo y los positivos en negro; el $0$ va en negro con borde grueso.
  \item \textbf{Control.} Verifica que no haya valores repetidos y que las marcas del piso coincidan con las tarjetas.
\end{enumerate}

\subsection{Patrones explícitos (fabricación)}

\subsubsection{Patrón A — Las 20 tarjetas}

\begin{center}
\small
\begin{tabular}{@{}r r@{\hspace{1.4cm}} r r@{\hspace{1.4cm}} r r@{\hspace{1.4cm}} r r@{}}
\toprule
\textbf{\#} & \textbf{Valor} & \textbf{\#} & \textbf{Valor} & \textbf{\#} & \textbf{Valor} & \textbf{\#} & \textbf{Valor}\\
\midrule
1 & $-10$ & 6 & $-5$ & 11 & $0$ & 16 & $+5$\\
2 & $-9$ & 7 & $-4$ & 12 & $+1$ & 17 & $+6$\\
3 & $-8$ & 8 & $-3$ & 13 & $+2$ & 18 & $+7$\\
4 & $-7$ & 9 & $-2$ & 14 & $+3$ & 19 & $+8$\\
5 & $-6$ & 10 & $-1$ & 15 & $+4$ & 20 & $+9$\\
\bottomrule
\end{tabular}
\end{center}

\subsubsection{Patrón B — Disposición de la cinta}

\begin{center}
\small
\begin{tabular}{@{}l c@{}}
\toprule
\textbf{Zona} & \textbf{Marcas}\\
\midrule
Izquierda & $-10$ a $-1$ (diez marcas)\\
Centro & $0$ (una marca, con borde grueso)\\
Derecha & $+1$ a $+10$ (diez marcas)\\
\bottomrule
\end{tabular}
\end{center}

\subsubsection{Patrón C — Frases de comparación}

\begin{center}
\small
\begin{tabular}{@{}p{5.4cm} p{8.8cm}@{}}
\toprule
\textbf{Situación} & \textbf{Frase válida}\\
\midrule
Número positivo & «Un negativo cualquiera es menor y un positivo mayor es mayor».\\
Número negativo & «Un negativo más lejano del cero es menor; uno más cercano es mayor».\\
Cero & «Todo negativo es menor y todo positivo es mayor».\\
\bottomrule
\end{tabular}
\end{center}

\begin{cuadroadvertencia}
\textbf{Cómo fabricar otra cinta.} Usa un rango menor (por ejemplo, $-6$ a $+6$) para grupos que recién comienzan, o un rango mayor (hasta $\pm15$) para reforzar. Mantén siempre el $0$ al centro y la misma separación entre marcas.
\end{cuadroadvertencia}

\subsection{Configuración inicial}

\begin{enumerate}[leftmargin=*]
  \item Extiende la cinta en el piso, lejos de mesas y sillas.
  \item Baraja las 20 tarjetas y déjalas boca abajo.
  \item Reparte tres tarjetas a cada jugador.
  \item Entrega una hoja de puntos por jugador.
  \item Designa un juez (docente o estudiante rotativo) que valide ubicaciones y comparaciones.
\end{enumerate}

% ============================================================
\clearpage
\section{Juego 3 — El ascensor de dados: sube y baja}
% ============================================================

\subsection{Ficha rápida}

\begin{center}
\small
\begin{tabular}{@{}p{3cm} p{12.2cm}@{}}
\toprule
\textbf{Jugadores} & 2 a 4\\
\textbf{Duración} & 15--20 minutos\\
\textbf{Objetivo} & Sumar y restar enteros pequeños moviendo una ficha en una tira numérica de $-10$ a $+10$.\\
\textbf{Subtemas} & Suma y resta de una operación, desplazamiento en la recta.\\
\textbf{Qué se evalúa} & Que anuncie la operación y el resultado antes de mover la ficha.\\
\textbf{Nivel} & Iniciación.\\
\bottomrule
\end{tabular}
\end{center}

\subsection{La idea en una frase}

\begin{cuadroextra}
  «El dado no siempre sube: tiene cargas positivas y negativas. Tu ficha es un ascensor que debe llegar exactamente al piso $+10$.»
\end{cuadroextra}

\subsection{Cómo se juega en cinco pasos}

\begin{enumerate}[leftmargin=*]
  \item Todos colocan su ficha en el $0$.
  \item En tu turno lanzas el dado de cargas.
  \item Anuncias la operación y el resultado en voz alta.
  \item Mueves la ficha a la derecha si la cara es positiva y a la izquierda si es negativa.
  \item Gana quien llegue exactamente a $+10$.
\end{enumerate}

\subsection{Reglas}

\begin{enumerate}[leftmargin=*]
  \item \textbf{Un dado de cargas.} Las caras son $+1,\ +2,\ +3,\ -1,\ -2,\ -3$. Si no tienes un dado en blanco, usa la variante de la moneda y el dado común (Patrón D).
  \item \textbf{Partida en cero.} Todas las fichas comienzan en la casilla $0$.
  \item \textbf{Una tirada por turno.} Se lanza una sola vez; no se repite la tirada aunque el resultado no guste.
  \item \textbf{Anuncio antes de mover.} El jugador dice la operación completa, por ejemplo: «$3+(-2)=+1$».
  \item \textbf{Dirección según la cara.} Cara positiva: avanza a la derecha esa cantidad. Cara negativa: retrocede a la izquierda esa cantidad. El resultado de la suma solo indica la nueva posición.
  \item \textbf{Suelo.} Nadie baja de $-10$: si el desplazamiento llevaría la ficha más abajo, se queda en $-10$.
  \item \textbf{Techo con rebote.} Si el desplazamiento supera $+10$, el exceso se cuenta hacia atrás. Desde $+9$ con $+3$: $+10$, $+9$, $+8$; la ficha queda en $+8$.
  \item \textbf{Error de cálculo.} Si el resultado anunciado es incorrecto, el jugador no mueve y pasa el turno. Los rivales pueden comprobar en la tira.
  \item \textbf{Victoria exacta.} Gana quien caiga exactamente en $+10$. Si se acaba el tiempo, gana quien esté más cerca de $+10$; desempate: quien esté más alto.
\end{enumerate}

\subsection{Turno paso a paso}

\begin{enumerate}[leftmargin=*]
  \item \textbf{Lanza} el dado de cargas.
  \item \textbf{Lee} la cara que salió.
  \item \textbf{Di} la operación completa con el resultado.
  \item \textbf{Mueve} la ficha contando las casillas en voz alta.
  \item \textbf{Comprueba} si llegaste exactamente a $+10$.
\end{enumerate}

\subsection{Final y puntuación}

\begin{center}
\small
\begin{tabular}{@{}p{7.4cm} p{6.6cm}@{}}
\toprule
\textbf{Situación} & \textbf{Resultado}\\
\midrule
Cae exactamente en $+10$ & Gana la partida.\\
Se acaba el tiempo & Gana quien esté más cerca de $+10$.\\
Empate en distancia & Gana quien esté más alto.\\
\bottomrule
\end{tabular}
\end{center}

\subsection{Ejemplo resuelto}

La ficha está en $+7$ y sale $-2$.

\begin{cuadrosolucion}
\begin{enumerate}[leftmargin=*]
  \item Anuncia: «$7+(-2)=+5$».
  \item Como la cara es $-2$, retrocede dos casillas: $+6$, $+5$.
  \item La ficha queda en $+5$.
\end{enumerate}
Otro turno: la ficha está en $+9$ y sale $+3$.
\begin{enumerate}[leftmargin=*]
  \item Anuncia: «$9+3=+12$».
  \item Se pasa de $+10$ por dos casillas: rebota a $+9$ y $+8$.
  \item La ficha queda en $+8$.
\end{enumerate}
\end{cuadrosolucion}

\subsection{Materiales}

\begin{center}
\small
\begin{tabular}{@{}p{5.0cm} c p{7.4cm}@{}}
\toprule
\textbf{Material} & \textbf{Cantidad} & \textbf{Uso}\\
\midrule
Dado & 1 & Caras de carga\\
Cinta adhesiva o tirro & 1 rollo & Cubrir y rotular las caras\\
Cartulina o cartón & 1 pliego & Tira numérica\\
Fichas o tapas & 1 por jugador & Marcadores de posición\\
Marcadores & 2 & Rotular la tira y el dado\\
\bottomrule
\end{tabular}
\end{center}

\subsection{Construcción paso a paso}

\begin{enumerate}[leftmargin=*]
  \item \textbf{Dado.} Cubre las seis caras de un dado con tirro y escribe $+1,\ +2,\ +3,\ -1,\ -2,\ -3$.
  \item \textbf{Tira.} Recorta una tira de $60\text{ cm}\times 8\text{ cm}$ y divídela en $21$ casillas iguales.
  \item \textbf{Numeración.} Escribe de izquierda a derecha los valores $-10$ a $+10$.
  \item \textbf{Marcas.} Pinta el $0$ de verde («inicio») y el $+10$ de rojo («meta»).
  \item \textbf{Control.} Verifica que haya exactamente $21$ casillas y que el $0$ quede en el centro.
\end{enumerate}

\subsection{Patrones explícitos (fabricación)}

\subsubsection{Patrón A — Caras del dado de cargas}

\begin{center}
\small
\begin{tabular}{@{}l c c c c c c@{}}
\toprule
\textbf{Dado} & \textbf{Cara 1} & \textbf{Cara 2} & \textbf{Cara 3} & \textbf{Cara 4} & \textbf{Cara 5} & \textbf{Cara 6}\\
\midrule
Cargas & $+1$ & $+2$ & $+3$ & $-1$ & $-2$ & $-3$\\
\bottomrule
\end{tabular}
\end{center}

\subsubsection{Patrón B — La tira numérica}

\begin{center}
\small
\begin{tabular}{@{}l c@{}}
\toprule
\textbf{Tramo} & \textbf{Casillas}\\
\midrule
Izquierda & $-10,\ -9,\ -8,\ -7,\ -6,\ -5,\ -4,\ -3,\ -2,\ -1$\\
Centro & $0$\\
Derecha & $+1,\ +2,\ +3,\ +4,\ +5,\ +6,\ +7,\ +8,\ +9,\ +10$\\
\bottomrule
\end{tabular}
\end{center}

\subsubsection{Patrón C — Regla de movimiento}

\begin{center}
\small
\begin{tabular}{@{}p{5.4cm} p{8.8cm}@{}}
\toprule
\textbf{Cara del dado} & \textbf{Movimiento}\\
\midrule
Positiva & Avanza a la derecha esa cantidad.\\
Negativa & Retrocede a la izquierda esa cantidad.\\
Cero & No se mueve.\\
La ficha llegaría bajo $-10$ & Se queda en $-10$.\\
La ficha llegaría sobre $+10$ & Rebota: el exceso se cuenta hacia atrás.\\
\bottomrule
\end{tabular}
\end{center}

\subsubsection{Patrón D — Variante sin dado modificado}

Si no hay dado en blanco, se usa un dado común y una moneda:

\begin{enumerate}[leftmargin=*]
  \item Se lanza el dado y la moneda a la vez.
  \item Si la moneda sale \textbf{cara}, se suma el valor del dado; si sale \textbf{sello}, se resta.
  \item Se anuncia la operación, se determina la dirección y se mueve como en el juego normal.
\end{enumerate}

\begin{cuadroadvertencia}
\textbf{Cómo fabricar otra tira.} Si el grupo necesita un recorrido más corto, usa casillas de $-6$ a $+6$ y una meta en $+6$. Mantén siempre el $0$ al centro y una meta exacta.
\end{cuadroadvertencia}

\subsection{Configuración inicial}

\begin{enumerate}[leftmargin=*]
  \item Coloca la tira en el centro de la mesa.
  \item Cada jugador pone su ficha en la casilla $0$.
  \item Deja el dado de cargas al centro.
  \item Acuerda el orden de turno y el tiempo máximo (10 minutos).
\end{enumerate}

% ============================================================
\clearpage
\section{Juego 4 — Pesca de opuestos: el par perfecto}
% ============================================================

\subsection{Ficha rápida}

\begin{center}
\small
\begin{tabular}{@{}p{3cm} p{12.2cm}@{}}
\toprule
\textbf{Jugadores} & 2 a 4\\
\textbf{Duración} & 15--20 minutos\\
\textbf{Objetivo} & Reunir parejas de números opuestos pidiendo cartas a los compañeros.\\
\textbf{Subtemas} & Números opuestos y suma de opuestos igual a cero.\\
\textbf{Qué se evalúa} & Que nombre el opuesto correcto y explique por qué la pareja suma cero.\\
\textbf{Nivel} & Iniciación.\\
\bottomrule
\end{tabular}
\end{center}

\subsection{La idea en una frase}

\begin{cuadroextra}
  «Cada carta busca su reflejo: $+3$ necesita su $-3$. Pide, pesca y forma parejas; gana quien reúna más espejos.»
\end{cuadroextra}

\subsection{Cómo se juega en cinco pasos}

\begin{enumerate}[leftmargin=*]
  \item Se reparten cinco cartas a cada jugador.
  \item En tu turno eliges una carta y pides su opuesto a un rival.
  \item Si la tiene, te la entrega y formas la pareja.
  \item Si no la tiene, dices «¡A pescar!» y robas una carta del pozo.
  \item Gana quien reúna más parejas de opuestos.
\end{enumerate}

\subsection{Reglas}

\begin{enumerate}[leftmargin=*]
  \item \textbf{Mazo.} Se usan dos copias de cada entero de $-8$ a $+8$, sin el $0$: 32 cartas en total.
  \item \textbf{Reparto.} Se entregan cinco cartas a cada jugador; el resto forma el pozo, boca abajo al centro.
  \item \textbf{Pedir el opuesto.} En tu turno eliges una carta de tu mano y pides su opuesto a un rival, nombrando las dos: «¿Tienes el opuesto de $+3$, es decir, $-3$?».
  \item \textbf{Carta propia.} Solo puedes pedir el opuesto de una carta que tengas en la mano.
  \item \textbf{Si el rival la tiene.} Debe entregar una copia de esa carta. Formas la pareja, la muestras, ganas $1$ punto y juegas de nuevo.
  \item \textbf{Si no la tiene.} Dice «¡A pescar!» y robas una carta del pozo. Si la carta robada completa una pareja con tu mano, la muestras y sumas $1$ punto; luego pasa el turno.
  \item \textbf{Pareja válida.} Solo vale el par formado por un número y su opuesto: $(+v,-v)$. No vale juntar dos cartas iguales.
  \item \textbf{Pozo vacío.} Si se acaba el pozo, se sigue pidiendo hasta que nadie pueda formar parejas.
  \item \textbf{Sin cartas.} Si te quedas sin cartas, robas una del pozo; si tampoco hay pozo, quedas fuera.
  \item \textbf{Victoria.} Gana quien tenga más parejas al final. Desempate: la pareja con mayor valor absoluto.
\end{enumerate}

\subsection{Turno paso a paso}

\begin{enumerate}[leftmargin=*]
  \item \textbf{Mira} tu mano y elige una carta.
  \item \textbf{Calcula} su opuesto.
  \item \textbf{Pide} a un rival: «¿Tienes el opuesto de…?».
  \item \textbf{Recibe} la carta o pesca del pozo.
  \item \textbf{Muestra} la pareja si la formaste y anota el punto.
  \item \textbf{Repite} el turno si acertaste; si no, pasa el turno.
\end{enumerate}

\subsection{Final y puntuación}

\begin{center}
\small
\begin{tabular}{@{}p{8.2cm} c@{}}
\toprule
\textbf{Logro} & \textbf{Puntos}\\
\midrule
Pareja de opuestos formada & $1$\\
Pareja de cartas iguales & No vale\\
\bottomrule
\end{tabular}
\end{center}

Gana quien tenga más parejas cuando se acaben las posibilidades de pesca. Desempate: la pareja con mayor valor absoluto.

\subsection{Ejemplo resuelto}

El jugador tiene $+3$ en la mano y pide a un rival el opuesto.

\begin{cuadrosolucion}
\begin{enumerate}[leftmargin=*]
  \item Pregunta: «¿Tienes el opuesto de $+3$, es decir, $-3$?».
  \item El rival tiene $-3$ y se la entrega.
  \item El jugador muestra la pareja $(+3,-3)$ y explica: «$3+(-3)=0$, son opuestos».
  \item Gana $1$ punto y juega otra vez.
\end{enumerate}
Si el rival no tuviera $-3$:
\begin{enumerate}[leftmargin=*]
  \item Dice «¡A pescar!».
  \item El jugador roba una carta. Si sale $-3$, forma la pareja y suma; si no, pasa el turno.
\end{enumerate}
\end{cuadrosolucion}

\subsection{Materiales}

\begin{center}
\small
\begin{tabular}{@{}p{5.0cm} c p{7.4cm}@{}}
\toprule
\textbf{Material} & \textbf{Cantidad} & \textbf{Uso}\\
\midrule
Cartulina & 3 pliegos & 32 cartas\\
Marcadores negro y rojo & 1 de cada uno & Positivos y negativos\\
Tijera, regla y lápiz & 1 de cada uno & Trazado y corte\\
Bolsa o caja pequeña & 1 & Pozo de cartas\\
Hoja de puntos & 1 por jugador & Llevar el puntaje\\
\bottomrule
\end{tabular}
\end{center}

\subsection{Construcción paso a paso}

\begin{enumerate}[leftmargin=*]
  \item \textbf{Corte.} Divide la cartulina en 32 rectángulos de $6\text{ cm}\times 9\text{ cm}$.
  \item \textbf{Rotulado.} Escribe dos copias de cada entero de $-8$ a $+8$, sin el $0$ (Patrón A).
  \item \textbf{Colores.} Negativos en rojo y positivos en negro.
  \item \textbf{Reverso.} Dibuja un pez en el dorso de las 32 cartas para que no se distingan.
  \item \textbf{Control.} Cuenta 16 valores distintos, dos copias de cada uno, y confirma que no hay ningún $0$.
\end{enumerate}

\subsection{Patrones explícitos (fabricación)}

\subsubsection{Patrón A — Composición del mazo}

\begin{center}
\small
\begin{tabular}{@{}r c c@{\hspace{1.2cm}} r c c@{}}
\toprule
\textbf{Valor} & \textbf{Copias $+v$} & \textbf{Copias $-v$} & \textbf{Valor} & \textbf{Copias $+v$} & \textbf{Copias $-v$}\\
\midrule
1 & 2 & 2 & 5 & 2 & 2\\
2 & 2 & 2 & 6 & 2 & 2\\
3 & 2 & 2 & 7 & 2 & 2\\
4 & 2 & 2 & 8 & 2 & 2\\
\bottomrule
\end{tabular}
\end{center}

Total: $8\times 4=32$ cartas. Cada valor absoluto permite formar dos parejas de opuestos.

\subsubsection{Patrón B — Parejas válidas}

\begin{center}
\small
\begin{tabular}{@{}c c c@{}}
\toprule
\textbf{Pareja} & \textbf{Suma} & \textbf{Verificación}\\
\midrule
$(+1,-1)$ & $0$ & $1+(-1)=0$\\
$(+2,-2)$ & $0$ & $2+(-2)=0$\\
$(+3,-3)$ & $0$ & $3+(-3)=0$\\
$(+4,-4)$ & $0$ & $4+(-4)=0$\\
$(+5,-5)$ & $0$ & $5+(-5)=0$\\
$(+6,-6)$ & $0$ & $6+(-6)=0$\\
$(+7,-7)$ & $0$ & $7+(-7)=0$\\
$(+8,-8)$ & $0$ & $8+(-8)=0$\\
\bottomrule
\end{tabular}
\end{center}

\subsubsection{Patrón C — Guion de la pregunta}

\begin{center}
\small
\begin{tabular}{@{}p{6.2cm} p{8cm}@{}}
\toprule
\textbf{Momento} & \textbf{Frase exacta}\\
\midrule
Al pedir & «¿Tienes el opuesto de $+v$, es decir, $-v$?»\\
Si tiene la carta & «Sí, aquí está; forma la pareja».\\
Si no tiene la carta & «¡A pescar!».\\
Al mostrar la pareja & «Son opuestos porque suman $0$».\\
\bottomrule
\end{tabular}
\end{center}

\begin{cuadroadvertencia}
\textbf{Cómo fabricar otro mazo.} Usa dos copias de cada entero de $-12$ a $+12$ (sin el $0$) para un juego más largo, o de $-6$ a $+6$ para uno más corto. Mantén siempre dos copias por valor y excluye el $0$, que no tiene un opuesto distinto.
\end{cuadroadvertencia}

\subsection{Configuración inicial}

\begin{enumerate}[leftmargin=*]
  \item Baraja las 32 cartas.
  \item Reparte cinco cartas a cada jugador.
  \item Coloca el resto en el pozo, boca abajo al centro.
  \item Entrega una hoja de puntos por jugador.
  \item Acuerda quién empieza (por ejemplo, quien tenga la carta de mayor valor absoluto).
\end{enumerate}

% ============================================================
\clearpage
\section{Juego 5 — Tres en raya de enteros: conquista el resultado}
% ============================================================

\subsection{Ficha rápida}

\begin{center}
\small
\begin{tabular}{@{}p{3cm} p{12.2cm}@{}}
\toprule
\textbf{Jugadores} & 2 jugadores o 2 equipos\\
\textbf{Duración} & 15--20 minutos\\
\textbf{Objetivo} & Resolver sumas y restas y marcar el resultado en una cuadrícula $3\times 3$.\\
\textbf{Subtemas} & Suma y resta de enteros de una operación.\\
\textbf{Qué se evalúa} & Que resuelva en voz alta y verifique el resultado antes de marcar.\\
\textbf{Nivel} & Iniciación.\\
\bottomrule
\end{tabular}
\end{center}

\subsection{La idea en una frase}

\begin{cuadroextra}
  «El tablero no se juega con X y O: se juega con resultados. Si resuelves la operación, conquistas la casilla; si te equivocas, la pierdes.»
\end{cuadroextra}

\subsection{Cómo se juega en cinco pasos}

\begin{enumerate}[leftmargin=*]
  \item Se dibuja una cuadrícula $3\times 3$ con nueve resultados.
  \item Cada jugador elige un color o un montón de fichas.
  \item En tu turno sacas una tarjeta y resuelves la operación.
  \item Si el resultado está libre en el tablero, lo marcas con tu color.
  \item Gana quien alinee tres marcas propias.
\end{enumerate}

\subsection{Reglas}

\begin{enumerate}[leftmargin=*]
  \item \textbf{Tablero fijo.} La cuadrícula tiene los nueve resultados del Patrón A y no se cambia durante la partida.
  \item \textbf{Mazo.} Se usan 27 tarjetas: tres operaciones para cada uno de los nueve resultados (Patrón B).
  \item \textbf{Turnos alternos.} El Jugador 1 empieza y luego se alterna. En equipos, cada integrante juega en su turno.
  \item \textbf{Resolver en voz alta.} Se saca la tarjeta superior, se lee la operación y se resuelve en voz alta.
  \item \textbf{Marcar.} Si el resultado es correcto y la casilla está libre, se marca con el color propio (o se coloca una ficha).
  \item \textbf{Sin marca.} Si el resultado es incorrecto, la casilla está ocupada o el resultado no está en el tablero, no se marca y pasa el turno.
  \item \textbf{Auditoría.} El rival puede pedir la explicación del cálculo. Si el jugador no puede justificarlo, se retira la marca y pasa el turno.
  \item \textbf{No borrar.} Una marca puesta no se mueve ni se borra, salvo por la auditoría del punto anterior.
  \item \textbf{Victoria.} Gana quien alinee tres marcas propias en fila, columna o diagonal.
  \item \textbf{Mazo agotado.} Si se acaban las 27 tarjetas y nadie alineó tres, gana quien tenga más marcas. Si hay empate, se roban tarjetas extra hasta que alguien marque.
\end{enumerate}

\subsection{Turno paso a paso}

\begin{enumerate}[leftmargin=*]
  \item \textbf{Saca} la tarjeta superior del mazo.
  \item \textbf{Lee} la operación en voz alta.
  \item \textbf{Resuelve} y di el resultado.
  \item \textbf{Marca} la casilla si está libre.
  \item \textbf{Comprueba} si lograste tres en raya.
  \item \textbf{Pasa} el turno.
\end{enumerate}

\subsection{Final y puntuación}

\begin{center}
\small
\begin{tabular}{@{}p{7.6cm} p{6.4cm}@{}}
\toprule
\textbf{Situación} & \textbf{Resultado}\\
\midrule
Tres marcas en fila, columna o diagonal & Gana la partida.\\
Mazo agotado sin tres en raya & Gana quien tenga más marcas.\\
Empate en marcas & Se roban tarjetas extra hasta que alguien marque.\\
\bottomrule
\end{tabular}
\end{center}

\subsection{Ejemplo resuelto}

Sale la tarjeta $-3+(-3)$.

\begin{cuadrosolucion}
\begin{enumerate}[leftmargin=*]
  \item El jugador resuelve: «$-3+(-3)=-6$».
  \item Mira el tablero: el $-6$ está en la esquina superior izquierda y está libre.
  \item Coloca su ficha o marca el $-6$ con su color.
  \item Verifica si tiene tres en raya; si no, pasa el turno.
\end{enumerate}
Turno siguiente: sale $9+(-4)$.
\begin{enumerate}[leftmargin=*]
  \item Resuelve: «$9+(-4)=+5$».
  \item El $+5$ está libre en la esquina superior derecha: lo marca.
\end{enumerate}
\end{cuadrosolucion}

\subsection{Materiales}

\begin{center}
\small
\begin{tabular}{@{}p{5.0cm} c p{7.4cm}@{}}
\toprule
\textbf{Material} & \textbf{Cantidad} & \textbf{Uso}\\
\midrule
Cartulina o cartón & 1 pliego & Tablero $3\times 3$\\
Cartulina & 2 pliegos & 27 tarjetas\\
Fichas de dos colores & 10 de cada color & Marcar casillas\\
Marcadores & 2 & Rotular tablero y tarjetas\\
Tijera, regla y lápiz & 1 de cada uno & Trazado y corte\\
\bottomrule
\end{tabular}
\end{center}

\subsection{Construcción paso a paso}

\begin{enumerate}[leftmargin=*]
  \item \textbf{Tablero.} Recorta un cuadrado de $18\text{ cm}\times 18\text{ cm}$ y divídelo en una cuadrícula $3\times 3$ de $6\text{ cm}$ por lado.
  \item \textbf{Valores.} Copia los nueve resultados del Patrón A, uno por casilla.
  \item \textbf{Tarjetas.} Recorta 27 rectángulos de $6\text{ cm}\times 9\text{ cm}$ y copia las operaciones del Patrón B.
  \item \textbf{Reverso.} Escribe la respuesta al dorso de cada tarjeta para la auditoría.
  \item \textbf{Control.} Verifica que cada resultado aparezca exactamente en tres tarjetas y que todas las operaciones sean correctas.
\end{enumerate}

\subsection{Patrones explícitos (fabricación)}

\subsubsection{Patrón A — El tablero $3\times 3$}

\begin{center}
\small
\begin{tabular}{|c|c|c|}
\hline
$-6$ & $+1$ & $+5$\\\hline
$-2$ & $0$ & $+7$\\\hline
$+3$ & $-4$ & $+9$\\\hline
\end{tabular}
\end{center}

\subsubsection{Patrón B — Las 27 tarjetas}

\begin{center}
\footnotesize
\begin{tabular}{@{}r l c@{\hspace{0.8cm}} r l c@{\hspace{0.8cm}} r l c@{}}
\toprule
\textbf{\#} & \textbf{Operación} & \textbf{Result.} & \textbf{\#} & \textbf{Operación} & \textbf{Result.} & \textbf{\#} & \textbf{Operación} & \textbf{Result.}\\
\midrule
1 & $-3+(-3)$ & $-6$ & 10 & $5+(-7)$ & $-2$ & 19 & $-5+8$ & $+3$\\
2 & $-9+3$ & $-6$ & 11 & $-8+6$ & $-2$ & 20 & $9-6$ & $+3$\\
3 & $2-8$ & $-6$ & 12 & $3-5$ & $-2$ & 21 & $-7+10$ & $+3$\\
4 & $-4+5$ & $+1$ & 13 & $-4+4$ & $0$ & 22 & $2+(-6)$ & $-4$\\
5 & $7-6$ & $+1$ & 14 & $6-6$ & $0$ & 23 & $-9+5$ & $-4$\\
6 & $-8+9$ & $+1$ & 15 & $-9+9$ & $0$ & 24 & $1-5$ & $-4$\\
7 & $9+(-4)$ & $+5$ & 16 & $12+(-5)$ & $+7$ & 25 & $4+5$ & $+9$\\
8 & $-2+7$ & $+5$ & 17 & $-1+8$ & $+7$ & 26 & $12-3$ & $+9$\\
9 & $11-6$ & $+5$ & 18 & $10-3$ & $+7$ & 27 & $-2+11$ & $+9$\\
\bottomrule
\end{tabular}
\end{center}

\subsubsection{Patrón C — Clave por resultado}

\begin{center}
\small
\begin{tabular}{@{}c c@{}}
\toprule
\textbf{Resultado} & \textbf{Tarjetas que lo producen}\\
\midrule
$-6$ & 1, 2, 3\\
$+1$ & 4, 5, 6\\
$+5$ & 7, 8, 9\\
$-2$ & 10, 11, 12\\
$0$ & 13, 14, 15\\
$+7$ & 16, 17, 18\\
$+3$ & 19, 20, 21\\
$-4$ & 22, 23, 24\\
$+9$ & 25, 26, 27\\
\bottomrule
\end{tabular}
\end{center}

\begin{cuadroadvertencia}
\textbf{Cómo fabricar otro tablero.} Elige nueve resultados distintos entre $-12$ y $+12$ y escribe tres operaciones para cada uno. Comprueba que \textbf{todos} los resultados del mazo estén en el tablero y que ninguna operación dé un resultado ajeno. Si cambias un valor del tablero, debes cambiar sus tres tarjetas.
\end{cuadroadvertencia}

\subsection{Configuración inicial}

\begin{enumerate}[leftmargin=*]
  \item Coloca el tablero en el centro.
  \item Baraja las 27 tarjetas y déjalas boca abajo.
  \item Entrega cinco fichas de un color a cada jugador o equipo.
  \item Acuerda quién empieza y deja una hoja de borrador para las comprobaciones.
\end{enumerate}

% ============================================================
\clearpage
\section{Juego 6 — Guerra de pisos: el ascensor y los sótanos}
% ============================================================

\subsection{Ficha rápida}

\begin{center}
\small
\begin{tabular}{@{}p{3cm} p{12.2cm}@{}}
\toprule
\textbf{Jugadores} & 2 a 4 (con 2 es el más fluido)\\
\textbf{Duración} & 15--20 minutos\\
\textbf{Objetivo} & Comparar enteros y justificar el orden usando la recta numérica vertical del edificio.\\
\textbf{Subtemas} & Relación de orden ($<$, $>$, $=$) y ubicación en la recta numérica.\\
\textbf{Qué se evalúa} & Que el estudiante diga la relación correcta y la explique con la posición en la recta.\\
\textbf{Nivel} & Entrada (más sencillo que el Juego 11).\\
\bottomrule
\end{tabular}
\end{center}

\subsection{La idea en una frase}

\begin{cuadroextra}
  «Cada carta es un piso del edificio: los negativos son sótanos y los positivos son pisos altos. En cada ronda se voltea una carta y gana el piso más alto. Lo importante no es la suerte: es saber leer el edificio.»
\end{cuadroextra}

\subsection{Cómo se juega en cinco pasos}

\begin{enumerate}[leftmargin=*]
  \item Se barajan las 40 cartas y se reparten en partes iguales, \textbf{boca abajo}.
  \item Todos voltean a la vez la carta superior de su montón.
  \item Quien tiene la carta \textbf{mayor} dice la relación en voz alta y la justifica.
  \item Ese jugador se lleva las cartas de la ronda y las coloca al fondo de su montón.
  \item Gana quien se quede con todas las cartas (o quien tenga más al terminar el tiempo).
\end{enumerate}

\subsection{Reglas}

\textbf{Idea clave:} cada ronda es una comparación sobre la recta vertical del edificio; el objetivo es que el ganador \emph{explique} por qué su carta es mayor.

\begin{enumerate}[leftmargin=*]
  \item \textbf{Reparto.} Con 40 cartas: 20 y 20 si juegan 2; 13 y 13 con 14 aparte si juegan 3; 10 cada uno si juegan 4.
  \item \textbf{Volteo simultáneo.} Todos voltean la carta superior de su montón a la vez. Nadie puede mirar las cartas antes de la señal.
  \item \textbf{Quién gana la ronda.} Gana la carta de mayor valor. Ese jugador debe anunciar la relación completa y justificarla: «$+3>-5$, porque $+3$ está más arriba en la recta».
  \item \textbf{Si el ganador no justifica.} La ronda se anula: las cartas van al montón de apartadas y nadie las gana.
  \item \textbf{Si el ganador se equivoca.} Si compara mal o justifica mal, las cartas de la ronda las recoge el jugador con la \textbf{segunda carta más alta}. Si solo quedan dos jugadores, las recoge el otro.
  \item \textbf{Empate (guerra de pisos).} Si dos o más cartas son iguales, cada jugador implicado coloca una carta boca abajo y voltea otra. Quien gane esa segunda comparación se lleva \textbf{todas} las cartas acumuladas.
  \item \textbf{Límite de guerras.} No se pueden encadenar más de 3 guerras seguidas. En la cuarta, las cartas acumuladas se apartan y se inicia una ronda nueva.
  \item \textbf{El $0$ es la planta baja.} Se compara como cualquier entero: es menor que todos los positivos y mayor que todos los negativos.
  \item \textbf{Cartas apartadas.} Las cartas de una ronda anulada o de una guerra cortada van a un montón aparte y no vuelven a la partida.
  \item \textbf{Eliminación.} Si un jugador se queda sin cartas al final de una ronda, queda eliminado. Gana el último jugador que conserve cartas o, si se acaba el tiempo, quien tenga más.
  \item \textbf{Comprobación.} Cualquier jugador puede pedir que se revise una comparación. Si el ganador se equivocó, se aplica la regla 5.
\end{enumerate}

\subsection{Turno paso a paso}

\begin{enumerate}[leftmargin=*]
  \item \textbf{Voltea} tu carta superior cuando todos estén listos.
  \item \textbf{Mira} las cartas de los demás y localiza la mayor.
  \item \textbf{Compara} mentalmente con la recta: más arriba, mayor.
  \item \textbf{Anuncia} la relación y justifícala si tu carta es la mayor.
  \item \textbf{Recoge} las cartas si ganaste, o entrégalas si te equivocaste.
  \item \textbf{Coloca} las cartas ganadas al fondo de tu montón.
  \item \textbf{Comprueba} si alguien se quedó sin cartas.
\end{enumerate}

\subsection{Final y puntuación}

\begin{center}
\small
\begin{tabular}{@{}p{7.4cm} p{6.6cm}@{}}
\toprule
\textbf{Final} & \textbf{Quién gana}\\
\midrule
Un jugador reúne las 40 cartas & Ese jugador\\
Se acaba el tiempo acordado (10 min) & Quien tenga más cartas\\
Empate en cantidad & Quien tenga la carta de mayor valor\\
\bottomrule
\end{tabular}
\end{center}

\subsection{Ejemplo resuelto}

Voltean tres jugadores: $+3$, $-5$ y $0$.

\begin{cuadrosolucion}
\begin{enumerate}[leftmargin=*]
  \item Orden en la recta: $-5<0<+3$.
  \item Gana el jugador del $+3$ y debe decir: «$+3>0$ y $+3>-5$, porque $+3$ está más arriba».
  \item Si dijera «$+3<-5$», el jugador del $0$ (el siguiente mayor) se lleva las tres cartas.
  \item El ganador coloca las tres cartas al fondo de su montón.
\end{enumerate}
\end{cuadrosolucion}

\subsection{Materiales}

\begin{center}
\small
\begin{tabular}{@{}p{5.0cm} c p{7.4cm}@{}}
\toprule
\textbf{Material} & \textbf{Cantidad} & \textbf{Uso}\\
\midrule
Cartulina o cartón & 2 pliegos & 40 cartas\\
Marcador negro y rojo & 1 de cada uno & Positivos y negativos\\
Regla, lápiz y tijera & 1 de cada uno & Trazado y corte\\
Reloj o cronómetro & 1 & Controlar los 10 minutos\\
\bottomrule
\end{tabular}
\end{center}

\subsection{Construcción paso a paso}

\begin{enumerate}[leftmargin=*]
  \item \textbf{Corte.} Divide la cartulina en \textbf{40 rectángulos de $6\text{ cm}\times 9\text{ cm}$}.
  \item \textbf{Rotulado.} Escribe en el centro de cada carta un entero, en negro si es positivo o cero y en rojo si es negativo. Usa el Patrón A.
  \item \textbf{Reverso.} Dibuja en el dorso un pequeño edificio con sótanos y pisos, para que todas las cartas se vean iguales por detrás.
  \item \textbf{Control.} Verifica que haya exactamente \textbf{dos cartas de cada valor} y ninguna repetida de más.
\end{enumerate}

\subsection{Patrones explícitos (fabricación)}

\subsubsection{Patrón A — Composición del mazo (40 cartas)}

El mazo es \textbf{dos copias de cada entero desde $-10$ hasta $+9$}: 20 valores distintos, 40 cartas.

\begin{center}
\small
\begin{tabular}{@{}r r@{\hspace{1.4cm}} r r@{\hspace{1.4cm}} r r@{\hspace{1.4cm}} r r@{}}
\toprule
\textbf{\#} & \textbf{Valor} & \textbf{\#} & \textbf{Valor} & \textbf{\#} & \textbf{Valor} & \textbf{\#} & \textbf{Valor}\\
\midrule
1 & $-10$ & 6 & $-5$ & 11 & $0$ & 16 & $+5$\\
2 & $-9$ & 7 & $-4$ & 12 & $+1$ & 17 & $+6$\\
3 & $-8$ & 8 & $-3$ & 13 & $+2$ & 18 & $+7$\\
4 & $-7$ & 9 & $-2$ & 14 & $+3$ & 19 & $+8$\\
5 & $-6$ & 10 & $-1$ & 15 & $+4$ & 20 & $+9$\\
\bottomrule
\end{tabular}
\end{center}

\begin{cuadroadvertencia}
\textbf{Cómo fabricar otra baraja.} Elige 20 enteros consecutivos y haz dos copias de cada uno. Si quieres un juego más largo, usa 30 valores; si quieres uno más corto, usa 12.
\end{cuadroadvertencia}

\subsubsection{Patrón B — Las tres justificaciones válidas}

\begin{center}
\small
\begin{tabular}{@{}p{4.4cm} p{9.6cm}@{}}
\toprule
\textbf{Situación} & \textbf{Justificación que debe decir el estudiante}\\
\midrule
Ambos positivos & «El de mayor valor numérico está más arriba».\\
Un positivo y un negativo & «Todo positivo está por encima de todo negativo».\\
Ambos negativos & «El de mayor valor absoluto está más abajo, así que es el menor».\\
\bottomrule
\end{tabular}
\end{center}

\subsubsection{Patrón C — Regla de desempate}

\begin{center}
\small
\begin{tabular}{@{}p{5.4cm} p{8.6cm}@{}}
\toprule
\textbf{Ronda} & \textbf{Qué se hace}\\
\midrule
Empate simple & Una carta boca abajo y se voltea otra.\\
Empate en la guerra & Se repite una vez más.\\
Tercera guerra seguida & Las cartas se apartan y se inicia una ronda nueva.\\
\bottomrule
\end{tabular}
\end{center}

\subsection{Configuración inicial}

\begin{enumerate}[leftmargin=*]
  \item Baraja las 40 cartas.
  \item Reparte en partes iguales y coloca cada montón \textbf{boca abajo} frente a su jugador.
  \item Deja el montón de cartas apartadas en el centro.
  \item Prepara el cronómetro en 10 minutos.
\end{enumerate}

% ============================================================
\clearpage
\section{Juego 7 — Memoria de opuestos: el espejo del cero}
% ============================================================

\subsection{Ficha rápida}

\begin{center}
\small
\begin{tabular}{@{}p{3cm} p{12.2cm}@{}}
\toprule
\textbf{Jugadores} & 2 a 4\\
\textbf{Duración} & 15--20 minutos\\
\textbf{Objetivo} & Reconocer números opuestos e iguales valores absolutos, y explicar por qué son equivalentes.\\
\textbf{Subtemas} & Números opuestos (simétricos), valor absoluto $\abs{a}$, suma de opuestos igual a cero.\\
\textbf{Qué se evalúa} & Que el estudiante justifique la pareja: opuestos o mismo valor absoluto.\\
\textbf{Nivel} & Entrada (versión simple del radar del Juego 11).\\
\bottomrule
\end{tabular}
\end{center}

\subsection{La idea en una frase}

\begin{cuadroextra}
  «Cada pareja es un reflejo en el espejo del cero: dos números opuestos, o un número y su distancia al cero. Encuentra el reflejo y explícalo.»
\end{cuadroextra}

\subsection{Cómo se juega en cinco pasos}

\begin{enumerate}[leftmargin=*]
  \item Se colocan las 40 cartas \textbf{boca abajo} en una cuadrícula de $8\times 5$.
  \item En tu turno volteas \textbf{dos cartas}.
  \item Si forman una pareja válida, te las quedas, explicas por qué y sumas los puntos.
  \item Si no forman pareja, las vuelves a tapar y pasa el turno.
  \item Gana quien sume más puntos al retirarse las 20 parejas.
\end{enumerate}

\subsection{Reglas}

\textbf{Idea clave:} el juego es de memoria; las cartas siempre vuelven a su posición y la pareja solo vale si se explica.

\begin{enumerate}[leftmargin=*]
  \item \textbf{Composición del mazo.} 40 cartas: 24 de opuestos (12 pares $\pm1$ a $\pm12$) y 16 de valor absoluto (8 pares $\abs{-v}$ con $+v$).
  \item \textbf{Disposición.} Las 40 cartas se colocan boca abajo en una cuadrícula de 8 columnas por 5 filas.
  \item \textbf{Dos cartas por turno.} Se voltean exactamente dos cartas y se comparan.
  \item \textbf{Parejas válidas.} Hay dos tipos:
  \begin{itemize}
    \item \textbf{Opuestos:} $x$ y $-x$ (por ejemplo, $+7$ y $-7$). Se justifica con «son opuestos porque suman $0$».
    \item \textbf{Mismo valor absoluto:} una carta $\abs{-15}$ y otra $+15$. Se justifica con «$\abs{-15}=15$; es el mismo valor».
  \end{itemize}
  \item \textbf{Explicación obligatoria.} Si el jugador no puede explicar la pareja, las cartas se vuelven a tapar y pierde el turno, aunque las cartas fueran correctas.
  \item \textbf{Turno extra.} Quien forma una pareja válida y la explica retira las dos cartas y juega otra vez.
  \item \textbf{Puntos.} Pareja de opuestos: $1$ punto. Pareja de valor absoluto bien explicada: $2$ puntos.
  \item \textbf{Sin pistas.} No se puede señalar, soplar ni tocar la ubicación de las cartas. Si ocurre, la pareja se aparta y no suma.
  \item \textbf{Posiciones fijas.} Al final de cada turno, toda carta que no se retiró vuelve exactamente a su lugar.
  \item \textbf{Final.} La partida termina cuando se retiran las 20 parejas. Gana quien sume más puntos; desempate: más parejas de valor absoluto.
\end{enumerate}

\subsection{Turno paso a paso}

\begin{enumerate}[leftmargin=*]
  \item \textbf{Elige} dos cartas y voltéalas.
  \item \textbf{Comprueba} si son opuestos o si tienen el mismo valor absoluto.
  \item \textbf{Explica} en voz alta la razón antes de retirarlas.
  \item \textbf{Si es pareja:} retíralas, suma los puntos y repite el turno.
  \item \textbf{Si no es pareja:} tápalas en el mismo lugar.
  \item \textbf{Pasa} el turno al siguiente jugador.
\end{enumerate}

\subsection{Final y puntuación}

\begin{center}
\small
\begin{tabular}{@{}p{8.6cm} c@{}}
\toprule
\textbf{Concepto} & \textbf{Puntos}\\
\midrule
Pareja de opuestos & $1$\\
Pareja de valor absoluto (bien explicada) & $2$\\
Pareja mal explicada & $0$ (se vuelve a tapar)\\
\bottomrule
\end{tabular}
\end{center}

Gana quien sume más puntos cuando se acaban las 20 parejas. Desempate: más parejas de valor absoluto.

\subsection{Ejemplo resuelto}

El jugador voltea $\mathbf{-9}$ y $\mathbf{+9}$.

\begin{cuadrosolucion}
\begin{enumerate}[leftmargin=*]
  \item Comprueba: $-9+9=0$.
  \item Explica: «Son opuestos: están a la misma distancia del cero y suman $0$».
  \item Retira la pareja y gana $1$ punto. Como acertó, juega otra vez.
\end{enumerate}
En el siguiente turno voltea $\abs{-15}$ y $+15$:
\begin{enumerate}[leftmargin=*]
  \item Resuelve: $\abs{-15}=15$.
  \item Explica: «Es el mismo valor, porque el valor absoluto es la distancia al cero».
  \item Retira la pareja y gana $2$ puntos.
\end{enumerate}
\end{cuadrosolucion}

\subsection{Materiales}

\begin{center}
\small
\begin{tabular}{@{}p{5.0cm} c p{7.4cm}@{}}
\toprule
\textbf{Material} & \textbf{Cantidad} & \textbf{Uso}\\
\midrule
Cartulina blanca & 2 pliegos & 40 cartas\\
Cartulina de color (opcional) & 1 pliego & Cartas de valor absoluto\\
Marcadores negro y rojo & 1 de cada uno & Rotular\\
Regla, lápiz y tijera & 1 de cada uno & Trazado y corte\\
\bottomrule
\end{tabular}
\end{center}

\subsection{Construcción paso a paso}

\begin{enumerate}[leftmargin=*]
  \item \textbf{Corte.} Divide la cartulina en \textbf{40 rectángulos de $6\text{ cm}\times 9\text{ cm}$}.
  \item \textbf{Cartas de opuestos (24).} Escribe los pares $\pm1$ a $\pm12$ del Patrón A, en rojo los negativos y en negro los positivos.
  \item \textbf{Cartas de valor absoluto (16).} En cartulina de color escribe las 8 expresiones $\abs{-v}$ del Patrón B y sus 8 resultados $+v$.
  \item \textbf{Reverso.} Dibuja el mismo símbolo del espejo en el dorso de las 40 cartas.
\end{enumerate}

\subsection{Patrones explícitos (fabricación)}

\subsubsection{Patrón A — 12 pares de opuestos (24 cartas)}

\begin{center}
\small
\begin{tabular}{@{}c l@{\hspace{1.2cm}} c l@{\hspace{1.2cm}} c l@{}}
\toprule
\textbf{Par} & \textbf{Cartas} & \textbf{Par} & \textbf{Cartas} & \textbf{Par} & \textbf{Cartas}\\
\midrule
1 & $+1$ y $-1$ & 5 & $+5$ y $-5$ & 9 & $+9$ y $-9$\\
2 & $+2$ y $-2$ & 6 & $+6$ y $-6$ & 10 & $+10$ y $-10$\\
3 & $+3$ y $-3$ & 7 & $+7$ y $-7$ & 11 & $+11$ y $-11$\\
4 & $+4$ y $-4$ & 8 & $+8$ y $-8$ & 12 & $+12$ y $-12$\\
\bottomrule
\end{tabular}
\end{center}

\subsubsection{Patrón B — 8 pares de valor absoluto (16 cartas)}

\begin{center}
\small
\begin{tabular}{@{}r c c l@{}}
\toprule
\textbf{\#} & \textbf{Carta A} & \textbf{Carta B} & \textbf{Verificación}\\
\midrule
1 & $\abs{-15}$ & $+15$ & $\abs{-15}=15$\\
2 & $\abs{-18}$ & $+18$ & $\abs{-18}=18$\\
3 & $\abs{-20}$ & $+20$ & $\abs{-20}=20$\\
4 & $\abs{-25}$ & $+25$ & $\abs{-25}=25$\\
5 & $\abs{-30}$ & $+30$ & $\abs{-30}=30$\\
6 & $\abs{-40}$ & $+40$ & $\abs{-40}=40$\\
7 & $\abs{-50}$ & $+50$ & $\abs{-50}=50$\\
8 & $\abs{-100}$ & $+100$ & $\abs{-100}=100$\\
\bottomrule
\end{tabular}
\end{center}

\begin{cuadroadvertencia}
\textbf{Cómo fabricar nuevos pares.} Usa valores que no aparezcan en el Patrón A, para que no haya dos cartas iguales en el mazo. En el Patrón B, todos los valores son mayores que 12.
\end{cuadroadvertencia}

\subsubsection{Patrón C — Cuadrícula de colocación}

Las 40 cartas se colocan en \textbf{8 columnas por 5 filas}. Al final de cada turno, cada carta vuelve exactamente a su posición: así el juego es de memoria y no de azar.

\subsection{Configuración inicial}

\begin{enumerate}[leftmargin=*]
  \item Baraja las 40 cartas.
  \item Colócalas boca abajo en una cuadrícula de 8 columnas por 5 filas.
  \item Deja una hoja de puntos por jugador.
  \item Acuerda quién empieza.
\end{enumerate}

% ============================================================
\clearpage
\section{Juego 8 — La serpiente de los enteros: suma y resta en un salto}
% ============================================================

\subsection{Ficha rápida}

\begin{center}
\small
\begin{tabular}{@{}p{3cm} p{12.2cm}@{}}
\toprule
\textbf{Jugadores} & 2 a 4\\
\textbf{Duración} & 20--25 minutos\\
\textbf{Objetivo} & Resolver sumas y restas de enteros de una sola operación y usar el resultado para avanzar o retroceder en un tablero.\\
\textbf{Subtemas} & Adición y sustracción de enteros, signo del resultado, desplazamiento en la recta.\\
\textbf{Qué se evalúa} & Que el estudiante resuelva la operación y explique si avanza o retrocede.\\
\textbf{Nivel} & Entrada (una sola operación; el Juego 12 trabaja las propiedades).\\
\bottomrule
\end{tabular}
\end{center}

\subsection{La idea en una frase}

\begin{cuadroextra}
  «La serpiente es una recta numérica del 0 al 30. Cada carta te da un salto: si el resultado es positivo avanzas, si es negativo retrocedes, y si es cero te quedas quieto.»
\end{cuadroextra}

\subsection{Cómo se juega en cinco pasos}

\begin{enumerate}[leftmargin=*]
  \item Todos los jugadores colocan su ficha en la casilla \textbf{0}.
  \item En tu turno sacas una carta del mazo y resuelves la operación en voz alta.
  \item Avanzas tantas casillas como indique el resultado; si es negativo, retrocedes.
  \item Si el salto te pasa del 30, rebotas: el exceso se cuenta hacia atrás.
  \item Gana quien caiga \textbf{exactamente} en la casilla 30.
\end{enumerate}

\subsection{Reglas}

\textbf{Idea clave:} una carta, una operación y un movimiento. El tablero es una recta del $0$ al $30$ y hay que caer exactamente en la meta.

\begin{enumerate}[leftmargin=*]
  \item \textbf{Un salto por turno.} Se saca una sola carta y se resuelve una sola operación.
  \item \textbf{Cálculo en voz alta.} El jugador dice la operación y el resultado antes de mover: «$2-(-3)=+5$; avanzo 5».
  \item \textbf{Dirección del movimiento.} Resultado positivo: avanza. Resultado negativo: retrocede. Resultado cero: no se mueve.
  \item \textbf{Suelo en 0.} Nadie puede quedar por debajo de la casilla $0$: si el retroceso lo llevaría por debajo de $0$, se queda en $0$.
  \item \textbf{Techo en 30 con rebote.} Si el salto supera el $30$, se cuenta el exceso hacia atrás. Desde $28$ con $+5$: $29, 30, 29, 28, 27$; la ficha queda en $27$.
  \item \textbf{Error de cálculo.} Si el resultado es incorrecto, el jugador \textbf{no mueve} y pasa el turno. Los rivales pueden pedir la comprobación en la recta del tablero.
  \item \textbf{Mazo agotado.} Cuando se acaban las 24 cartas, se barajan otra vez y se continúa.
  \item \textbf{Victoria exacta.} Gana quien caiga \textbf{exactamente} en la casilla $30$.
  \item \textbf{Empate en la misma ronda.} Si dos jugadores llegan al $30$ en la misma ronda de turnos, gana quien haya usado menos cartas.
  \item \textbf{Tiempo agotado.} Si se acaba el tiempo acordado, gana quien esté más cerca del $30$; desempate: quien tenga el turno.
\end{enumerate}

\subsection{Turno paso a paso}

\begin{enumerate}[leftmargin=*]
  \item \textbf{Saca} la carta superior del mazo.
  \item \textbf{Resuelve} la operación y di el resultado.
  \item \textbf{Decide} si avanzas, retrocedes o te quedas.
  \item \textbf{Mueve} la ficha contando las casillas en voz alta.
  \item \textbf{Comprueba} si llegaste exactamente al $30$.
  \item \textbf{Pasa} el turno si no ganaste.
\end{enumerate}

\subsection{Final y puntuación}

Gana quien llegue \textbf{exactamente} a la casilla 30. Si se acaba el tiempo acordado, gana quien esté más cerca del 30 (y, en caso de empate, quien tenga el turno).

\subsection{Ejemplo resuelto}

El jugador está en la casilla 12 y saca la carta $2-(-3)$.

\begin{cuadrosolucion}
\begin{enumerate}[leftmargin=*]
  \item Resuelve: $2-(-3)=2+3=+5$.
  \item Explica: «Restar un negativo es sumar; avanzo 5».
  \item Cuenta: 13, 14, 15, 16, 17. Queda en la casilla 17.
\end{enumerate}
Otro ejemplo con retroceso: está en 8 y saca $-6+2$.
\begin{enumerate}[leftmargin=*]
  \item Resuelve: $-6+2=-4$.
  \item Explica: «El resultado es negativo; retrocedo 4».
  \item Cuenta: 7, 6, 5, 4. Queda en la casilla 4.
\end{enumerate}
\end{cuadrosolucion}

\subsection{Materiales}

\begin{center}
\small
\begin{tabular}{@{}p{5.0cm} c p{7.4cm}@{}}
\toprule
\textbf{Material} & \textbf{Cantidad} & \textbf{Uso}\\
\midrule
Cartulina o cartón & 1 pliego grande & Tablero de 30 casillas\\
Cartulina & 1 pliego & 24 cartas\\
Fichas o tapas & 1 por jugador & Marcadores de posición\\
Marcadores negro y rojo & 1 de cada uno & Rotular\\
Regla, lápiz y tijera & 1 de cada uno & Trazado y corte\\
\bottomrule
\end{tabular}
\end{center}

\subsection{Construcción paso a paso}

\begin{enumerate}[leftmargin=*]
  \item \textbf{Tablero.} Dibuja una serpiente o un camino de \textbf{30 casillas numeradas del 0 al 30}, con curvas amplias para que quepan las fichas.
  \item \textbf{Casillas especiales.} Marca en color el 0 («inicio») y el 30 («meta»).
  \item \textbf{Cartas.} Recorta 24 tarjetas de $6\text{ cm}\times 9\text{ cm}$ y copia las operaciones del Patrón A.
  \item \textbf{Control.} Comprueba que ninguna operación dé un resultado mayor que 8 ni menor que $-6$.
\end{enumerate}

\subsection{Patrones explícitos (fabricación)}

\subsubsection{Patrón A — Las 24 cartas (mazo completo)}

\begin{center}
\footnotesize
\begin{tabular}{@{}r l r c@{\hspace{0.7cm}} r l r c@{}}
\toprule
\textbf{\#} & \textbf{Operación} & \textbf{Result.} & \textbf{Tipo} & \textbf{\#} & \textbf{Operación} & \textbf{Result.} & \textbf{Tipo}\\
\midrule
1 & $3+(-8)$ & $-5$ & Suma & 13 & $-4-(-3)$ & $-1$ & Resta\\
2 & $-6+2$ & $-4$ & Suma & 14 & $6-9$ & $-3$ & Resta\\
3 & $7+(-7)$ & $0$ & Suma & 15 & $0-6$ & $-6$ & Resta\\
4 & $5+(-2)$ & $+3$ & Suma & 16 & $-1-5$ & $-6$ & Resta\\
5 & $-2+7$ & $+5$ & Suma & 17 & $2-(-3)$ & $+5$ & Resta\\
6 & $-5+3$ & $-2$ & Suma & 18 & $-8-(-4)$ & $-4$ & Resta\\
7 & $8+(-5)$ & $+3$ & Suma & 19 & $3-(-4)$ & $+7$ & Resta\\
8 & $-6+8$ & $+2$ & Suma & 20 & $1-(-1)$ & $+2$ & Resta\\
9 & $4+(-4)$ & $0$ & Suma & 21 & $4-6$ & $-2$ & Resta\\
10 & $2+5$ & $+7$ & Suma & 22 & $-7-(-3)$ & $-4$ & Resta\\
11 & $-9+4$ & $-5$ & Suma & 23 & $6-2$ & $+4$ & Resta\\
12 & $1+6$ & $+7$ & Suma & 24 & $-1-(-1)$ & $0$ & Resta\\
\bottomrule
\end{tabular}
\end{center}

\textbf{Composición:} 12 sumas y 12 restas; resultados entre $-6$ y $+7$; tres cartas con resultado $0$.

\subsubsection{Patrón B — Regla del salto}

\begin{center}
\small
\begin{tabular}{@{}p{5.2cm} p{9cm}@{}}
\toprule
\textbf{Resultado} & \textbf{Movimiento}\\
\midrule
Positivo & Avanza esa cantidad de casillas.\\
Negativo & Retrocede esa cantidad de casillas (sin bajar de 0).\\
Cero & No se mueve.\\
Se pasa del 30 & Rebote: el exceso se cuenta hacia atrás.\\
\bottomrule
\end{tabular}
\end{center}

\begin{cuadroadvertencia}
\textbf{Cómo fabricar nuevas cartas.} Elige dos enteros entre $-9$ y $+9$ y una operación (suma o resta). Comprueba que el resultado quede entre $-6$ y $+8$ para que el tablero de 30 casillas siga siendo jugable.
\end{cuadroadvertencia}

\subsection{Configuración inicial}

\begin{enumerate}[leftmargin=*]
  \item Extiende el tablero en el centro de la mesa.
  \item Cada jugador coloca su ficha en la casilla \textbf{0}.
  \item Baraja las 24 cartas y déjalas boca abajo junto al tablero.
  \item Acuerda el orden de turno y deja una hoja de borrador por jugador.
\end{enumerate}

% ============================================================
\clearpage
\section{Juego 9 — Bingo de productos: la fábrica de signos}
% ============================================================

\subsection{Ficha rápida}

\begin{center}
\small
\begin{tabular}{@{}p{3cm} p{12.2cm}@{}}
\toprule
\textbf{Jugadores} & 2 a 4\\
\textbf{Duración} & 20--25 minutos\\
\textbf{Objetivo} & Calcular productos y cocientes exactos de enteros y aplicar la regla de los signos.\\
\textbf{Subtemas} & Multiplicación y división exacta en $\Z$, regla de los signos, divisor distinto de cero.\\
\textbf{Qué se evalúa} & Que el estudiante anuncie el resultado y justifique el signo.\\
\textbf{Nivel} & Entrada (una sola operación; el Juego 13 incluye distributiva y estrategia).\\
\bottomrule
\end{tabular}
\end{center}

\subsection{La idea en una frase}

\begin{cuadroextra}
  «Dos cartas deciden la operación: una da el primer factor y la otra el segundo. Si el resultado está en tu cartón, lo marcas; el primero en completar una línea gana.»
\end{cuadroextra}

\subsection{Cómo se juega en cinco pasos}

\begin{enumerate}[leftmargin=*]
  \item Cada jugador recibe un \textbf{cartón de $4\times 4$} con 16 resultados.
  \item Se barajan los dos mazos: \textbf{mazo A} y \textbf{mazo B}, con los enteros de $-6$ a $+6$ sin el cero.
  \item En cada ronda se voltea una carta de cada mazo y una moneda decide: \textbf{cara = multiplicar}, \textbf{sello = dividir}.
  \item Todos calculan la operación; quien tenga ese resultado en su cartón lo marca.
  \item Gana quien complete primero una fila, una columna o una diagonal.
\end{enumerate}

\subsection{Reglas}

\textbf{Idea clave:} cada ronda produce un solo resultado; si está en tu cartón, lo marcas, pero siempre anunciando la operación y el signo.

\begin{enumerate}[leftmargin=*]
  \item \textbf{Un cartón por jugador.} Cada jugador recibe un cartón $4\times 4$ con 16 resultados y 16 fichas para marcar.
  \item \textbf{Dos mazos.} El mazo A y el mazo B tienen una carta de cada entero de $-6$ a $+6$, sin el $0$ (12 cartas cada uno).
  \item \textbf{Una ronda, dos cartas.} En cada ronda se voltea una carta de cada mazo y se lanza una moneda: \textbf{cara = multiplicar}, \textbf{sello = dividir}.
  \item \textbf{Orden de la división.} Se divide la carta del mazo A entre la carta del mazo B, en el orden en que salieron.
  \item \textbf{División exacta.} Si el cociente no es exacto, \textbf{se multiplica obligatoriamente}. No se puede dividir entre $0$ porque los mazos no tienen ceros.
  \item \textbf{Canto obligatorio.} El jugador que marca debe decir la operación y el signo: «$-4\cdot 3=-12$; menos por más da menos». Si no lo dice, no puede marcar.
  \item \textbf{Marcado.} Se marca con una ficha o una X. Si el resultado no está en el cartón, no se marca nada.
  \item \textbf{Una casilla por ronda.} Aunque el resultado aparezca dos veces en el cartón, solo se marca una casilla por ronda.
  \item \textbf{Verificación.} Cualquier jugador puede pedir la comprobación. Si el que marcó se equivocó, retira la ficha.
  \item \textbf{Mazos agotados.} Se barajan de nuevo y se continúa.
  \item \textbf{Victoria.} Gana quien complete primero una fila, una columna o una diagonal de 4 casillas marcadas. Empate: gana quien tenga más casillas marcadas.
\end{enumerate}

\subsection{Turno paso a paso}

\begin{enumerate}[leftmargin=*]
  \item \textbf{Voltea} una carta del mazo A y una del mazo B.
  \item \textbf{Lanza} la moneda para decidir la operación.
  \item \textbf{Comprueba} que la división sea exacta; si no, multiplica.
  \item \textbf{Di} el resultado y justifica el signo.
  \item \textbf{Marca} una casilla si el resultado está en tu cartón.
  \item \textbf{Comprueba} si completaste una línea.
  \item \textbf{Pasa} el turno y rota quién voltea las cartas.
\end{enumerate}

\subsection{Final y puntuación}

Gana quien complete primero \textbf{una fila, una columna o una diagonal} de 4 casillas marcadas. En caso de empate, gana quien tenga más casillas marcadas en total.

\subsection{Ejemplo resuelto}

Se voltean las cartas $A=-4$ y $B=3$, y la moneda sale \textbf{cara (multiplicar)}.

\begin{cuadrosolucion}
\begin{enumerate}[leftmargin=*]
  \item Operación: $(-4)\cdot 3$.
  \item Signo: menos por más da menos.
  \item Resultado: $-12$.
  \item Si el cartón tiene el $-12$, el jugador lo marca y dice: «$-4$ por $3$ es $-12$».
\end{enumerate}
Si la moneda hubiera salido \textbf{sello (dividir)}: $(-4)\div 3$ no es exacta, así que se multiplica obligatoriamente y el resultado vuelve a ser $-12$.
\end{cuadrosolucion}

\subsection{Materiales}

\begin{center}
\small
\begin{tabular}{@{}p{5.0cm} c p{7.4cm}@{}}
\toprule
\textbf{Material} & \textbf{Cantidad} & \textbf{Uso}\\
\midrule
Cartulina & 3 pliegos & 4 cartones y los mazos A y B\\
Fichas, tapas o botones & 20 por jugador & Marcar casillas\\
Moneda & 1 & Elegir multiplicar o dividir\\
Marcadores & 2 & Rotular cartones y cartas\\
Regla, lápiz y tijera & 1 de cada uno & Trazado y corte\\
\bottomrule
\end{tabular}
\end{center}

\subsection{Construcción paso a paso}

\begin{enumerate}[leftmargin=*]
  \item \textbf{Cartones.} Recorta 4 cartones de $16\text{ cm}\times 16\text{ cm}$ y divide cada uno en una cuadrícula $4\times 4$ de $4\text{ cm}$ por lado.
  \item \textbf{Valores.} Copia en cada cartón los 16 valores del Patrón B.
  \item \textbf{Mazos.} Recorta 24 cartas de $5\text{ cm}\times 7\text{ cm}$: 12 para el mazo A y 12 para el mazo B, con los valores del Patrón A.
  \item \textbf{Reverso.} Usa un color para el mazo A y otro para el mazo B.
\end{enumerate}

\subsection{Patrones explícitos (fabricación)}

\subsubsection{Patrón A — Los dos mazos (12 cartas cada uno)}

\begin{center}
\small
\begin{tabular}{@{}l c c c c c c c c c c c c@{}}
\toprule
\textbf{Mazo A} & $-6$ & $-5$ & $-4$ & $-3$ & $-2$ & $-1$ & $+1$ & $+2$ & $+3$ & $+4$ & $+5$ & $+6$\\
\textbf{Mazo B} & $-6$ & $-5$ & $-4$ & $-3$ & $-2$ & $-1$ & $+1$ & $+2$ & $+3$ & $+4$ & $+5$ & $+6$\\
\bottomrule
\end{tabular}
\end{center}

Los dos mazos son \textbf{idénticos}: cada uno tiene una carta de cada entero de $-6$ a $+6$, sin el cero.

\subsubsection{Patrón B — Los 4 cartones}

Cada cartón es una cuadrícula $4\times 4$ y se lee por filas.

\begin{center}
\begin{minipage}{0.46\linewidth}\centering\small
  \textbf{Cartón 1}\\[2pt]
  \begin{tabular}{|c|c|c|c|}\hline
    $-16$ & $-4$ & $-12$ & $-2$\\\hline
    $-5$ & $-1$ & $+16$ & $+10$\\\hline
    $+5$ & $+12$ & $-9$ & $-15$\\\hline
    $+4$ & $+2$ & $+9$ & $-10$\\\hline
  \end{tabular}
\end{minipage}\hfill
\begin{minipage}{0.46\linewidth}\centering\small
  \textbf{Cartón 2}\\[2pt]
  \begin{tabular}{|c|c|c|c|}\hline
    $-15$ & $+1$ & $+12$ & $+10$\\\hline
    $+16$ & $+5$ & $-4$ & $-3$\\\hline
    $-20$ & $-9$ & $+8$ & $-6$\\\hline
    $-2$ & $-10$ & $+2$ & $-16$\\\hline
  \end{tabular}
\end{minipage}
\end{center}

\vspace{6pt}

\begin{center}
\begin{minipage}{0.46\linewidth}\centering\small
  \textbf{Cartón 3}\\[2pt]
  \begin{tabular}{|c|c|c|c|}\hline
    $-16$ & $-3$ & $+2$ & $+10$\\\hline
    $+3$ & $-4$ & $-8$ & $+1$\\\hline
    $-25$ & $-12$ & $-1$ & $+6$\\\hline
    $+4$ & $-5$ & $+9$ & $+8$\\\hline
  \end{tabular}
\end{minipage}\hfill
\begin{minipage}{0.46\linewidth}\centering\small
  \textbf{Cartón 4}\\[2pt]
  \begin{tabular}{|c|c|c|c|}\hline
    $+10$ & $+2$ & $-1$ & $+3$\\\hline
    $-10$ & $-20$ & $+9$ & $+1$\\\hline
    $-15$ & $-12$ & $-5$ & $-6$\\\hline
    $-8$ & $+12$ & $+5$ & $+16$\\\hline
  \end{tabular}
\end{minipage}
\end{center}

\subsubsection{Patrón C — Clave de operaciones (para el docente)}

Todos los valores de los cartones se pueden obtener con una operación de los dos mazos.

\begin{center}
\small
\begin{tabular}{@{}r l@{\hspace{1.1cm}} r l@{\hspace{1.1cm}} r l@{}}
\toprule
\textbf{Valor} & \textbf{Operación} & \textbf{Valor} & \textbf{Operación} & \textbf{Valor} & \textbf{Operación}\\
\midrule
$-25$ & $-5\cdot 5$ & $-5$ & $-5\cdot 1$ & $+3$ & $1\cdot 3$\\
$-20$ & $-5\cdot 4$ & $-4$ & $-2\cdot 2$ & $+4$ & $2\cdot 2$\\
$-16$ & $-4\cdot 4$ & $-3$ & $-3\cdot 1$ & $+5$ & $1\cdot 5$\\
$-15$ & $-5\cdot 3$ & $-2$ & $-2\cdot 1$ & $+6$ & $2\cdot 3$\\
$-12$ & $-4\cdot 3$ & $-1$ & $-1\cdot 1$ & $+8$ & $2\cdot 4$\\
$-10$ & $-5\cdot 2$ & $+1$ & $1\cdot 1$ & $+9$ & $3\cdot 3$\\
$-9$ & $-3\cdot 3$ & $+2$ & $1\cdot 2$ & $+10$ & $2\cdot 5$\\
$-8$ & $-4\cdot 2$ & $+12$ & $3\cdot 4$ & $+16$ & $4\cdot 4$\\
$-6$ & $-3\cdot 2$ & & & & \\
\bottomrule
\end{tabular}
\end{center}

\begin{cuadroadvertencia}
\textbf{Cómo fabricar otros cartones.} Elige 16 valores de la clave y repártelos en la cuadrícula sin repetir ninguno dentro del mismo cartón. Cualquier valor que no aparezca en la clave hace el cartón imposible de completar.
\end{cuadroadvertencia}

\subsection{Configuración inicial}

\begin{enumerate}[leftmargin=*]
  \item Entrega un cartón y 16 fichas a cada jugador.
  \item Baraja el mazo A y el mazo B por separado y colócalos boca abajo.
  \item Deja la moneda en el centro.
  \item Acuerda quién voltea las cartas en la primera ronda (se rota en cada ronda).
\end{enumerate}

% ============================================================
\clearpage
\section{Juego 10 — La escalera de dos peldaños: sube calculando}
% ============================================================

\subsection{Ficha rápida}

\begin{center}
\small
\begin{tabular}{@{}p{3cm} p{12.2cm}@{}}
\toprule
\textbf{Jugadores} & 2 a 4\\
\textbf{Duración} & 20--25 minutos\\
\textbf{Objetivo} & Resolver operaciones combinadas de \textbf{dos pasos} sin signos de agrupación, respetando la jerarquía.\\
\textbf{Subtemas} & Adición, sustracción, multiplicación y división exacta; prioridad de la multiplicación y la división.\\
\textbf{Qué se evalúa} & Que el estudiante resuelva primero el producto o cociente y después la suma o resta.\\
\textbf{Nivel} & Entrada (dos operaciones; el Juego 14 llega a cuatro operaciones con auditoría).\\
\bottomrule
\end{tabular}
\end{center}

\subsection{La idea en una frase}

\begin{cuadroextra}
  «Cada jugador tiene una escalera de 8 peldaños. Cada peldaño es una operación de dos pasos: si la resuelves bien, subes; si te equivocas, te quedas donde estás.»
\end{cuadroextra}

\subsection{Cómo se juega en cinco pasos}

\begin{enumerate}[leftmargin=*]
  \item Cada jugador recibe una \textbf{escalera} con 8 peldaños (una de las cuatro variantes).
  \item En tu turno resuelves el peldaño en el que estás.
  \item Si el resultado es correcto, \textbf{subes un peldaño}; si no, te quedas y pasa el turno.
  \item Los demás verifican la respuesta en la clave.
  \item Gana quien llegue primero al peldaño 8.
\end{enumerate}

\subsection{Reglas}

\textbf{Idea clave:} cada peldaño es una operación de dos pasos y el jugador debe mostrar el paso intermedio, no solo el resultado final.

\begin{enumerate}[leftmargin=*]
  \item \textbf{Una escalera por jugador.} Cada jugador recibe una escalera distinta (A, B, C o D) con 8 peldaños.
  \item \textbf{Un peldaño por turno.} Solo se resuelve el peldaño que corresponde; no se puede saltar ni adelantar.
  \item \textbf{Qué peldaño se resuelve.} La ficha marca el último peldaño superado. En el primer turno se resuelve el peldaño 1; después, siempre el siguiente al que marca la ficha.
  \item \textbf{Dos pasos, en orden.} Primero multiplicación o división; después adición o sustracción. El jugador debe decir el paso intermedio: «$4\cdot(-2)=-8$; luego $3-8=-5$».
  \item \textbf{Sin calculadora.} Se permite una hoja de borrador, pero no calculadora.
  \item \textbf{Respuesta completa.} Si el jugador dice solo el resultado final, sin el paso intermedio, la respuesta se considera incompleta y no sube.
  \item \textbf{Respuesta correcta.} Si el resultado coincide con la clave y el procedimiento se explicó, la ficha sube un peldaño. Si es incorrecto, no sube y pasa el turno.
  \item \textbf{Ayuda.} Si un jugador se traba, puede pedir una pista al docente, pero pierde el turno aunque después acierte. La pista se anota para el desempate.
  \item \textbf{Victoria.} Gana quien resuelva el peldaño 8 y suba la ficha a ese peldaño.
  \item \textbf{Empate.} Si dos jugadores llegan al peldaño 8 en la misma ronda, gana quien haya usado menos pistas; si persiste, quien haya llegado primero en el orden de turno.
\end{enumerate}

\subsection{Turno paso a paso}

\begin{enumerate}[leftmargin=*]
  \item \textbf{Lee} la expresión del peldaño que te corresponde.
  \item \textbf{Identifica} la operación de mayor jerarquía (producto o cociente).
  \item \textbf{Resuélvela} y di el resultado parcial.
  \item \textbf{Resuelve} la suma o resta final.
  \item \textbf{Di} el resultado completo en voz alta.
  \item \textbf{Sube} un peldaño si es correcto; si no, quédate.
  \item \textbf{Pasa} el turno.
\end{enumerate}

\subsection{Final y puntuación}

Gana quien llegue primero al \textbf{peldaño 8}. No hay puntos: la escalera mide avance.

\subsection{Ejemplo resuelto}

Peldaño: $3+4\cdot(-2)$.

\begin{cuadrosolucion}
\begin{enumerate}[leftmargin=*]
  \item Primero el producto: $4\cdot(-2)=-8$.
  \item Después la suma: $3+(-8)=-5$.
  \item El jugador dice: «$3+4\cdot(-2)=-5$» y sube un peldaño.
\end{enumerate}
\textbf{Error típico:} resolver de izquierda a derecha, $3+4=7$ y $7\cdot(-2)=-14$. Ese resultado no está en la clave.
\end{cuadrosolucion}

\subsection{Materiales}

\begin{center}
\small
\begin{tabular}{@{}p{5.0cm} c p{7.4cm}@{}}
\toprule
\textbf{Material} & \textbf{Cantidad} & \textbf{Uso}\\
\midrule
Cartulina & 2 pliegos & 4 escaleras (una por jugador)\\
Marcadores & 2 & Rotular peldaños\\
Regla, lápiz y tijera & 1 de cada uno & Trazado\\
Hoja de borrador & 1 por jugador & Cálculos\\
\bottomrule
\end{tabular}
\end{center}

\subsection{Construcción paso a paso}

\begin{enumerate}[leftmargin=*]
  \item \textbf{Escaleras.} Recorta 4 tiras de $10\text{ cm}\times 28\text{ cm}$ y dibuja 8 peldaños numerados de abajo hacia arriba.
  \item \textbf{Expresiones.} Escribe una expresión del Patrón A, B, C o D en cada peldaño (una variante por escalera).
  \item \textbf{Clave.} Copia la tabla de resultados del Patrón E en el reverso de cada escalera, para la verificación.
  \item \textbf{Forrado.} Si quieres reutilizarlas, pega las tiras sobre cartón y cúbrelas con cinta transparente.
\end{enumerate}

\subsection{Patrones explícitos (fabricación)}

\subsubsection{Patrón A — Escalera A (peldaños de abajo hacia arriba)}

\begin{center}
\small
\begin{tabular}{@{}r l r@{}}
\toprule
\textbf{Peldaño} & \textbf{Expresión} & \textbf{Resultado}\\
\midrule
1 & $-3+4\cdot 2$ & $+5$\\
2 & $5-3\cdot(-2)$ & $+11$\\
3 & $-2\cdot 3+8$ & $+2$\\
4 & $7+(-2)\cdot 4$ & $-1$\\
5 & $4\cdot(-1)-6$ & $-10$\\
6 & $-8+(-3)\cdot(-2)$ & $-2$\\
7 & $2\cdot 5-7$ & $+3$\\
8 & $-6-3\cdot(-2)$ & $0$\\
\bottomrule
\end{tabular}
\end{center}

\subsubsection{Patrón B — Escalera B}

\begin{center}
\small
\begin{tabular}{@{}r l r@{}}
\toprule
\textbf{Peldaño} & \textbf{Expresión} & \textbf{Resultado}\\
\midrule
1 & $-12\div 4+5$ & $+2$\\
2 & $8\div(-2)-3$ & $-7$\\
3 & $18\div(-3)+1$ & $-5$\\
4 & $-20\div 5-2$ & $-6$\\
5 & $6+(-15)\div 3$ & $+1$\\
6 & $-10\div(-2)-8$ & $-3$\\
7 & $3-(-24)\div 6$ & $+7$\\
8 & $-9+12\div(-4)$ & $-12$\\
\bottomrule
\end{tabular}
\end{center}

\subsubsection{Patrón C — Escalera C}

\begin{center}
\small
\begin{tabular}{@{}r l r@{}}
\toprule
\textbf{Peldaño} & \textbf{Expresión} & \textbf{Resultado}\\
\midrule
1 & $-5+3\cdot 4$ & $+7$\\
2 & $6-2\cdot(-2)$ & $+10$\\
3 & $-2\cdot(-3)+2$ & $+8$\\
4 & $9+(-4)\cdot 3$ & $-3$\\
5 & $-7-(-2)\cdot 5$ & $+3$\\
6 & $3\cdot(-3)+11$ & $+2$\\
7 & $-1+6\cdot 2$ & $+11$\\
8 & $8-(-4)\cdot(-2)$ & $0$\\
\bottomrule
\end{tabular}
\end{center}

\subsubsection{Patrón D — Escalera D}

\begin{center}
\small
\begin{tabular}{@{}r l r@{}}
\toprule
\textbf{Peldaño} & \textbf{Expresión} & \textbf{Resultado}\\
\midrule
1 & $15\div(-3)+7$ & $+2$\\
2 & $-8+20\div 4$ & $-3$\\
3 & $4-(-6)\div 2$ & $+7$\\
4 & $-2\cdot 4+11$ & $+3$\\
5 & $3+(-8)\cdot(-1)$ & $+11$\\
6 & $-14\div 2-(-3)$ & $-4$\\
7 & $6\cdot(-2)+11$ & $-1$\\
8 & $-5-9\div(-3)$ & $-2$\\
\bottomrule
\end{tabular}
\end{center}

\subsubsection{Patrón E — Regla de diseño}

\begin{center}
\small
\begin{tabular}{@{}p{4.6cm} p{9.4cm}@{}}
\toprule
\textbf{Elemento} & \textbf{Regla}\\
\midrule
Cantidad de operaciones & Exactamente dos: un producto o cociente y una suma o resta.\\
Rango de los resultados & Entre $-12$ y $+12$.\\
Divisiones & Siempre exactas y con divisor distinto de cero.\\
Dificultad & Los dos primeros peldaños son directos; del 5 al 8 incluyen más signos negativos.\\
Verificación & Cada escalera debe tener 8 resultados distintos.\\
\bottomrule
\end{tabular}
\end{center}

\begin{cuadroadvertencia}
\textbf{Cómo fabricar una escalera nueva.} Elige 8 resultados distintos entre $-12$ y $+12$; para cada uno construye una expresión de dos operaciones que lo produzca; comprueba el orden de las operaciones antes de escribirla en el peldaño.
\end{cuadroadvertencia}

\subsection{Configuración inicial}

\begin{enumerate}[leftmargin=*]
  \item Entrega una escalera distinta a cada jugador (A, B, C o D).
  \item Cada jugador coloca su ficha en el peldaño 0 («base»).
  \item Deja la clave de resultados boca abajo en el centro, para la verificación.
  \item Reparte una hoja de borrador y un lápiz a cada jugador.
\end{enumerate}

% ============================================================
\clearpage
\section{Juego 11 — Abismo y Cumbre: el radar del cero}
% ============================================================

\subsection{Ficha rápida}

\begin{center}
\small
\begin{tabular}{@{}p{3cm} p{12.2cm}@{}}
\toprule
\textbf{Jugadores} & 2 a 4\\
\textbf{Duración} & 30--40 minutos\\
\textbf{Objetivo} & Ubicar enteros en la recta numérica respetando el orden y reconocer opuestos y valor absoluto como distancia al cero.\\
\textbf{Subtemas} & Relación de orden ($<$, $>$, $=$), recta numérica, números opuestos, valor absoluto $\abs{a}$.\\
\textbf{Qué se evalúa} & Que el estudiante resuelva la expresión, anticipe el orden y explique por qué dos números son simétricos.\\
\textbf{Nivel} & Profundización.\\
\bottomrule
\end{tabular}
\end{center}

\subsection{La idea en una frase}

\begin{cuadroextra}
  «Somos operadores de un sonar: las tarjetas son coordenadas del fondo marino y la cuerda es la recta numérica. Si dos coordenadas están a la misma distancia del cero, el radar resuena.»
\end{cuadroextra}

\begin{center}
\begin{tikzpicture}[x=0.8cm]
  \draw[->, thick, mwprimario] (-7,0) -- (7,0);
  \foreach \x in {-6,-5,...,6} {
    \draw[mwprimario] (\x,0.12) -- (\x,-0.12);
  }
  \foreach \x in {-6,-4,-2,0,2,4,6} {
    \node[below, font=\scriptsize, color=mwgris] at (\x,-0.15) {$\x$};
  }
  \node[above, font=\scriptsize, color=mwprimario] at (0,0.35) {origen};
\end{tikzpicture}
\end{center}

\subsection{Cómo se juega en cinco pasos}

\begin{enumerate}[leftmargin=*]
  \item Se tensa la cuerda y se cuelga el $\mathbf{0}$ (origen) justo en el centro.
  \item Cada jugador recibe \textbf{4 tarjetas}, mirando solo la cara de la \textbf{expresión}.
  \item En tu turno eliges una tarjeta, \textbf{anuncias el hueco} donde la colocarás y resuelves la expresión.
  \item Volteas la tarjeta: si el valor cae dentro del hueco anunciado, la cuelgas y ganas \textbf{1 punto}; si su valor absoluto coincide con el de otra tarjeta ya colgada, ganas una \textbf{ficha de radar} (2 puntos).
  \item Robas hasta volver a tener 4 tarjetas. La partida termina cuando se agota el mazo y todos juegan su última mano.
\end{enumerate}

\subsection{Reglas}

\textbf{Idea clave:} se anuncia un hueco antes de conocer el valor; si el número cae dentro del hueco, la tarjeta se cuelga y, si comparte distancia al cero con otra, suena el radar.

\begin{enumerate}[leftmargin=*]
  \item \textbf{La cuerda es la recta.} Las tarjetas se cuelgan ordenadas de menor a mayor. No se repiten valores: toda tarjeta a la izquierda debe ser estrictamente menor y toda tarjeta a la derecha, estrictamente mayor.
  \item \textbf{Qué es un hueco.} Un hueco es el espacio entre dos tarjetas vecinas, o uno de los dos extremos de la cuerda.
  \item \textbf{Anuncio antes de resolver.} El jugador elige una tarjeta y anuncia el hueco \emph{antes} de voltearla: «La coloco entre $-25$ y $-12$».
  \item \textbf{Solo se mira la expresión.} Mientras elige, el jugador solo puede ver la cara A (la expresión). La cara B (el valor) se voltea después del anuncio.
  \item \textbf{Colocación válida.} Si el valor cae estrictamente dentro del hueco anunciado, la tarjeta se cuelga y el jugador gana $1$ punto. Si el valor es igual a un vecino, la colocación no es válida.
  \item \textbf{Error de estimación.} Si el valor no cae en el hueco anunciado, la tarjeta se descarta, no hay puntos y termina el turno. No se puede recolocar en otro hueco.
  \item \textbf{Resonancia de radar.} Después de colgar una tarjeta, si su valor absoluto coincide con el de otra tarjeta ya presente en la cuerda, se canta «¡Misma distancia al cero!», se explica y se toma $1$ ficha de radar ($2$ puntos). Cada tarjeta activa la resonancia una sola vez, al ser colgada.
  \item \textbf{Sondas de rango.} Las 6 tarjetas de sonda se resuelven al salir del mazo: el jugador gana $1$ punto por cada tarjeta ya colgada que cumpla la condición, aparta la sonda y roba otra tarjeta. Las sondas no se cuelgan en la cuerda.
  \item \textbf{Extremos.} También se puede anunciar «extremo izquierdo» o «extremo derecho»; la desigualdad se comprueba contra la única tarjeta vecina.
  \item \textbf{Final.} La partida termina cuando se agota el mazo y todos han jugado las tarjetas que les quedan en la mano.
  \item \textbf{Desempate.} Más puntos; si hay empate: (a) más fichas de radar; (b) más puntos de sondas; (c) más tarjetas propias en la cuerda.
\end{enumerate}

\subsection{Turno paso a paso}

\begin{enumerate}[leftmargin=*]
  \item \textbf{Mira} tus 4 tarjetas por la cara de la expresión.
  \item \textbf{Elige} una y anuncia el hueco exacto.
  \item \textbf{Resuelve} la expresión y di el valor.
  \item \textbf{Voltea} la tarjeta y comprueba el hueco.
  \item \textbf{Cuelga} la tarjeta si es válida; si no, descártala y pierde el turno.
  \item \textbf{Comprueba la resonancia} y toma la ficha de radar si corresponde.
  \item \textbf{Roba} hasta tener 4 tarjetas; si sale una sonda, resuélvela, apártala y sigue robando.
  \item \textbf{Pasa} el turno.
\end{enumerate}

\subsection{Final y puntuación}

La partida termina cuando \textbf{se agota el mazo de coordenadas y todos han jugado las tarjetas que les quedan en la mano}.

\begin{center}
\small
\begin{tabular}{@{}l c@{}}
\toprule
\textbf{Concepto} & \textbf{Puntos}\\
\midrule
Cada tarjeta propia correctamente colgada & $1$\\
Cada ficha de radar obtenida & $2$\\
Cada tarjeta que cumpla la condición de una sonda (al salir la sonda) & $1$\\
\bottomrule
\end{tabular}
\end{center}

Gana quien sume más puntos. Para partidas de 2 jugadores se pueden retirar al azar 8 coordenadas antes de empezar.

\subsection{Ejemplo resuelto}

En la cuerda están: $-25,\ -12,\ 0,\ +8,\ +18$. El jugador tiene la tarjeta «el opuesto de $\abs{-8}$».

\begin{cuadrosolucion}
\begin{enumerate}[leftmargin=*]
  \item Anuncia: «La coloco entre el $-12$ y el $0$».
  \item Resuelve: $\abs{-8}=8$; el opuesto de $8$ es $-8$.
  \item Comprueba: $-12<-8<0$. (\checkmark) Cuelga la tarjeta y gana $1$ punto.
  \item Radar: $\abs{-8}=8$ y en la cuerda está el $+8$, luego $\abs{-8}=\abs{+8}$. Toma una ficha de radar y explica: «$-8$ y $+8$ están a la misma distancia del cero». Gana $2$ puntos más.
  \item Roba una tarjeta.
\end{enumerate}
\textbf{Total de la jugada: $3$ puntos.}
\end{cuadrosolucion}

\subsection{Materiales}

\begin{center}
\small
\begin{tabular}{@{}p{5.0cm} c p{7.4cm}@{}}
\toprule
\textbf{Material} & \textbf{Cantidad} & \textbf{Uso}\\
\midrule
Cartulina blanca & 2 pliegos & 30 tarjetas (origen + coordenadas)\\
Cartulina de color & 1 pliego & 6 tarjetas de sonda\\
Pabilo, estambre grueso o cuerda fina & $1{,}5$ m & Recta numérica\\
Clips metálicos o ganchos pequeños & 36 & Sujetar tarjetas\\
Tapas plásticas marcadas con valor absoluto & 15 & Fichas de radar\\
Marcadores & negro, azul, rojo & Positivos y negativos\\
Regla de $30$ cm, tijera punta roma, lápiz & 1 cada uno & Trazado y corte\\
\bottomrule
\end{tabular}
\end{center}

\subsection{Construcción paso a paso}

\begin{enumerate}[leftmargin=*]
  \item \textbf{Mazo base.} Divide la cartulina blanca en \textbf{36 rectángulos de $6\text{ cm}\times 9\text{ cm}$} y recórtalos.
  \item \textbf{Origen.} Pinta un borde grueso negro en una tarjeta y escribe \textbf{0} en ambas caras: «Nivel del Mar / Origen».
  \item \textbf{Coordenadas (29 tarjetas).} En la \textbf{cara A} va la expresión; en la \textbf{cara B}, el valor entero simplificado. Dobla el borde superior $1$ cm para sujetarla con el clip.
  \item \textbf{Sondas (6 tarjetas).} En cartulina de color, recorta 6 tarjetas de $6\text{ cm}\times 9\text{ cm}$ y copia las condiciones del Patrón D.
  \item \textbf{Cuerda.} Corta $1{,}5$ m de pabilo, haz un lazo en cada extremo y fíjalo a dos sillas o al borde de la mesa.
\end{enumerate}

\subsection{Patrones explícitos (fabricación)}

\subsubsection{Patrón A — Tipos de expresión de la cara A}

\begin{center}
\small
\begin{tabular}{@{}c l l c@{}}
\toprule
\textbf{Código} & \textbf{Patrón} & \textbf{Ejemplo} & \textbf{Resultado}\\
\midrule
P1 & $-\abs{a}$ con $a>0$ & $-\abs{7}$ & $-7$\\
P2 & $\abs{-a}$ & $\abs{-9}$ & $+9$\\
P3 & $-\abs{-a}$ & $-\abs{-6}$ & $-6$\\
P4 & $\abs{a}$ & $\abs{5}$ & $+5$\\
P5 & «el opuesto de $x$» & opuesto de $(-14)$ & $+14$\\
P6 & $-(-a)$ & $-(-8)$ & $+8$\\
P7 & $-(+a)$ & $-(+5)$ & $-5$\\
\bottomrule
\end{tabular}
\end{center}

\subsubsection{Patrón B — Las 29 tarjetas de coordenadas}

\begin{longtable}{@{}r l c c@{}}
\toprule
\textbf{\#} & \textbf{Cara A (expresión)} & \textbf{Cara B} & \textbf{Patrón}\\
\midrule
\endhead
\bottomrule
\endfoot
1 & $-\abs{25}$ & $-25$ & P1\\
2 & $\abs{25}$ & $+25$ & P4\\
3 & opuesto de $\abs{20}$ & $-20$ & P5\\
4 & $\abs{-20}$ & $+20$ & P2\\
5 & $-\abs{-18}$ & $-18$ & P3\\
6 & $\abs{-18}$ & $+18$ & P2\\
7 & $-(+16)$ & $-16$ & P7\\
8 & $-(-16)$ & $+16$ & P6\\
9 & opuesto de $(-14)$ & $+14$ & P5\\
10 & $-\abs{14}$ & $-14$ & P1\\
11 & opuesto de $(+12)$ & $-12$ & P5\\
12 & $\abs{-12}$ & $+12$ & P2\\
13 & $-\abs{-9}$ & $-9$ & P3\\
14 & $\abs{-9}$ & $+9$ & P2\\
15 & opuesto de $\abs{-8}$ & $-8$ & P5\\
16 & $-(-8)$ & $+8$ & P6\\
17 & $-\abs{7}$ & $-7$ & P1\\
18 & opuesto de $(-7)$ & $+7$ & P5\\
19 & $-\abs{-6}$ & $-6$ & P3\\
20 & $\abs{-6}$ & $+6$ & P2\\
21 & $-(+5)$ & $-5$ & P7\\
22 & $\abs{5}$ & $+5$ & P4\\
23 & opuesto de $(-4)$ & $+4$ & P5\\
24 & $-(+4)$ & $-4$ & P7\\
25 & $-\abs{3}$ & $-3$ & P1\\
26 & opuesto de $(-3)$ & $+3$ & P5\\
27 & $-\abs{-2}$ & $-2$ & P3\\
28 & opuesto de $(-2)$ & $+2$ & P5\\
29 & $-\abs{30}$ & $-30$ & P1\\
\end{longtable}

\subsubsection{Patrón C — Pares simétricos del mazo}

El radar funciona porque el mazo contiene \textbf{14 pares de opuestos}. La tarjeta 29 ($-30$) no tiene pareja y es la «fosa abisal» de un solo uso.
\[
(\pm 25),\ (\pm 20),\ (\pm 18),\ (\pm 16),\ (\pm 14),\ (\pm 12),\ (\pm 9),\ (\pm 8),\ (\pm 7),\ (\pm 6),\ (\pm 5),\ (\pm 4),\ (\pm 3),\ (\pm 2).
\]

\subsubsection{Patrón D — Las 6 tarjetas de sonda}

\begin{center}
\small
\begin{tabular}{@{}c p{6.2cm} p{6.6cm}@{}}
\toprule
\textbf{\#} & \textbf{Condición de la sonda} & \textbf{Tarjetas del mazo que la cumplen}\\
\midrule
S1 & $\abs{x}\le 8$ y $x<0$ & $-8,-7,-6,-5,-4,-3,-2$ (7)\\
S2 & $\abs{x}\ge 12$ y $x>0$ & $+12,+14,+16,+18,+20,+25$ (6)\\
S3 & $x<0$ y $\abs{x}>10$ & $-30,-25,-20,-18,-16,-14,-12$ (7)\\
S4 & $x>0$ y $x<\abs{-9}$ & $+2,+3,+4,+5,+6,+7,+8$ (7)\\
S5 & $\abs{x}=\abs{-12}$ & $-12,+12$ (2)\\
S6 & $x<0$ y $\abs{x}\le 5$ & $-5,-4,-3,-2$ (4)\\
\bottomrule
\end{tabular}
\end{center}

\begin{cuadroadvertencia}
\textbf{Cómo fabricar nuevas sondas.} Elige una condición de signo y una de distancia; comprueba cuántas tarjetas del mazo la cumplen y escribe ese número al dorso como control. Si la cumplen más de 8 tarjetas, la sonda es demasiado fácil: cambia el umbral.
\end{cuadroadvertencia}

\subsection{Configuración inicial}

\begin{enumerate}[leftmargin=*]
  \item Tensa el pabilo horizontalmente a lo largo de la mesa.
  \item Cuelga la tarjeta del \textbf{0} exactamente en el centro.
  \item Baraja las 29 coordenadas y las 6 sondas juntas, con la cara A hacia arriba.
  \item Reparte \textbf{4 tarjetas de coordenadas} a cada jugador (si sale una sonda al repartir, se resuelve y se roba otra).
  \item Coloca las 15 tapas (fichas de radar) en el centro.
\end{enumerate}


% ============================================================
\clearpage
\section{Juego 12 — Laboratorio de Antimateria: el reactor}
% ============================================================

\subsection{Ficha rápida}

\begin{center}
\small
\begin{tabular}{@{}p{3cm} p{12.2cm}@{}}
\toprule
\textbf{Jugadores} & 2 a 4\\
\textbf{Duración} & 25--35 minutos\\
\textbf{Objetivo} & Comprender la adición y la sustracción de enteros como inyección y extracción de cargas, y usar las propiedades del $0$, del opuesto y de la asociatividad.\\
\textbf{Subtemas} & Adición y sustracción en $\Z$, $-(-a)=+a$, elemento neutro, inverso aditivo, conmutativa y asociativa.\\
\textbf{Qué se evalúa} & Que el estudiante traduzca la operación a un movimiento físico de tapas y explique el resultado.\\
\textbf{Nivel} & Profundización.\\
\bottomrule
\end{tabular}
\end{center}

\subsection{La idea en una frase}

\begin{cuadroextra}
  «Tu vaso es un reactor: las tapas azules son protones ($+1$) y las rojas son antiprotones ($-1$). Inyectar cambia la carga; extraer también. La meta es estabilizar el reactor en la carga exacta.»
\end{cuadroextra}

\subsection{Cómo se juega en cinco pasos}

\begin{enumerate}[leftmargin=*]
  \item Cada jugador recibe un vaso vacío y lo carga con 3 tapas al azar; anula los pares y anuncia su carga inicial.
  \item Se revela una \textbf{Carta de Meta} (por ejemplo, $-4$).
  \item En tu turno lanzas el \textbf{dado de operación} y el \textbf{dado de partícula}.
  \item Ejecutas el movimiento: inyectas, extraes o creas vacío; cancelas pares y anuncias la carga neta.
  \item Si tu carga coincide con la meta, reclamas la carta. Gana quien reúna \textbf{5 Cartas de Meta}.
\end{enumerate}

\subsection{Reglas}

\textbf{Idea clave:} cada operación es un movimiento físico de tapas; la carga del vaso es siempre azules menos rojas y los pares azul-rojo se aniquilan.

\begin{enumerate}[leftmargin=*]
  \item \textbf{Carga neta.} Carga $=$ (tapas azules) $-$ (tapas rojas). Un vaso vacío tiene carga $0$.
  \item \textbf{Aniquilación.} Un par azul-rojo vale $0$: se retira en cualquier momento sin cambiar la carga.
  \item \textbf{Inyectar es sumar.} Con $+$, se introducen las tapas que indique el dado de partícula: azules si es positivo, rojas si es negativo.
  \item \textbf{Extraer es restar.} Con $-$, se sacan del vaso las tapas indicadas. Extraer negativas (rojas) \emph{aumenta} la carga; extraer positivas (azules) la \emph{disminuye}.
  \item \textbf{Creación de vacío.} Si debes extraer $k$ tapas de un color y no tienes suficientes, añades primero $k$ pares nulos ($k$ azules y $k$ rojas) y luego extraes las $k$ tapas pedidas. No se puede decir «menos por menos da más» sin hacer la maniobra física.
  \item \textbf{Dado asociativo.} Si sale «Asociativa», se lanza dos veces el dado de partícula, se aplican ambos valores y se demuestra con las tapas que $(a+b)+c$ y $a+(b+c)$ dejan la misma carga neta.
  \item \textbf{Escudo neutro.} Cada propiedad declarada correctamente (neutro, opuesto o asociativa) da un escudo. Un escudo permite repetir una tirada desfavorable del dado de partícula, una vez por turno; se gasta al usarlo.
  \item \textbf{Meta.} Si al final del turno la carga coincide exactamente con la Carta de Meta, se reclama la carta y se revela otra. El reactor \textbf{no se vacía}: conserva su carga.
  \item \textbf{Mazo agotado.} Si se acaban las Cartas de Meta, se barajan las reclamadas y se continúa.
  \item \textbf{Victoria.} Gana el primero que reúna 5 Cartas de Meta. Empate: meta extra; gana quien la alcance primero.
\end{enumerate}

\subsection{Turno paso a paso}

\begin{enumerate}[leftmargin=*]
  \item \textbf{Lanza} el dado de operación y el dado de partícula.
  \item \textbf{Anuncia} la operación con tu carga actual: «Tengo $+2$ y salió extraer $-3$; es decir, $(+2)-(-3)$».
  \item \textbf{Ejecuta}: inyecta, extrae (con vacío si falta) o aplica la asociativa.
  \item \textbf{Aniquila} todos los pares nulos posibles.
  \item \textbf{Declara} la carga neta: «Quedan 5 azules; carga $+5$».
  \item \textbf{Comprueba la meta.} Si coincide, reclama la carta y revela otra.
  \item \textbf{Anota escudos} si declaraste una propiedad y pasa el turno.
\end{enumerate}

\subsection{Final y puntuación}

Gana el primero que reúna \textbf{5 Cartas de Meta de Estabilidad} («Jefe del Laboratorio»).

\begin{center}
\small
\begin{tabular}{@{}l l@{}}
\toprule
\textbf{Concepto} & \textbf{Efecto}\\
\midrule
Carta de Meta reclamada & $1$ punto de estabilidad\\
Escudo neutro & Permite repetir una tirada del dado de partícula\\
Empate & Meta extra; gana quien la alcance primero\\
\bottomrule
\end{tabular}
\end{center}

\subsection{Ejemplo resuelto}

El reactor tiene carga $-1$ (una tapa roja). La meta es $+2$.

\begin{cuadrosolucion}
\begin{enumerate}[leftmargin=*]
  \item Lanza: operación \textbf{extraer ($-$)} y partícula $\mathbf{-3}$.
  \item Anuncia: «Tengo $-1$ y debo extraer $-3$; es decir, $(-1)-(-3)$».
  \item No tiene 3 tapas rojas. \textbf{Crea vacío}: añade 2 pares nulos (2 azules y 2 rojas). Ahora hay 2 azules y 3 rojas.
  \item \textbf{Extrae} las 3 rojas. Quedan 2 azules.
  \item Carga neta: $+2$. Coincide con la meta: reclama la carta.
  \item Explica: $(-1)-(-3)=-1+3=+2$; al quitar cargas negativas el reactor sube.
\end{enumerate}
\end{cuadrosolucion}

\subsection{Materiales}

\begin{center}
\small
\begin{tabular}{@{}p{5.0cm} c p{7.4cm}@{}}
\toprule
\textbf{Material} & \textbf{Cantidad} & \textbf{Uso}\\
\midrule
Vasos plásticos transparentes o de cartón & 4 & Reactores\\
Tapas azules marcadas $+1$ & 40 & Protones\\
Tapas rojas marcadas $-1$ & 40 & Antiprotones\\
Dados & 2 & Operación y partícula\\
Etiquetas blancas o tirro & 1 rollo & Modificar las caras\\
Cartulina & 1 lámina & 20 Cartas de Meta\\
Marcadores & azul, rojo, negro & Rotular\\
Tijera, regla, pega & 1 cada uno & Armado\\
\bottomrule
\end{tabular}
\end{center}

\subsection{Construcción paso a paso}

\begin{enumerate}[leftmargin=*]
  \item \textbf{Partículas.} Escribe $+1$ en las 40 tapas azules y $-1$ en las 40 rojas. Si no hay tapas, recorta 80 círculos de cartulina de $3$ cm de diámetro.
  \item \textbf{Dado de operación.} Cubre las 6 caras de un dado con tirro y escribe: $+$ en 3 caras, $-$ en 2 caras y «Asociativa» en 1 cara.
  \item \textbf{Dado de partícula.} Cubre el segundo dado y escribe: $+1,\ -2,\ +3,\ -4,\ +5,\ 0$.
  \item \textbf{Cartas de Meta.} Recorta 20 rectángulos de $5\text{ cm}\times 7\text{ cm}$ y copia las metas del Patrón B.
  \item \textbf{Reactores.} Rotula cada vaso con tirro: «Reactor $\Z$».
\end{enumerate}

\subsection{Patrones explícitos (fabricación)}

\subsubsection{Patrón A — Caras de los dados}

\begin{center}
\small
\begin{tabular}{@{}l l@{}}
\toprule
\textbf{Dado} & \textbf{Caras}\\
\midrule
Operación & $+,\ +,\ +,\ -,\ -,\ \text{Asociativa}$\\
Partícula & $+1,\ -2,\ +3,\ -4,\ +5,\ 0$\\
\bottomrule
\end{tabular}
\end{center}

\textbf{Por qué estos valores:} el $0$ permite demostrar el elemento neutro; hay dos negativos y tres positivos para que la carga no se dispare siempre en un mismo sentido; ningún valor supera $5$ en magnitud para que la meta sea alcanzable en pocos turnos.

\subsubsection{Patrón B — Las 20 Cartas de Meta}

\begin{center}
\small
\begin{tabular}{@{}r r@{\hspace{1.6cm}} r r@{\hspace{1.6cm}} r r@{\hspace{1.6cm}} r r@{}}
\toprule
\textbf{\#} & \textbf{Meta} & \textbf{\#} & \textbf{Meta} & \textbf{\#} & \textbf{Meta} & \textbf{\#} & \textbf{Meta}\\
\midrule
1 & $-10$ & 6 & $-5$ & 11 & $0$ & 16 & $+5$\\
2 & $-9$  & 7 & $-4$ & 12 & $+1$ & 17 & $+6$\\
3 & $-8$  & 8 & $-3$ & 13 & $+2$ & 18 & $+7$\\
4 & $-7$  & 9 & $-2$ & 14 & $+3$ & 19 & $+8$\\
5 & $-6$  & 10 & $-1$ & 15 & $+4$ & 20 & $+9$\\
\bottomrule
\end{tabular}
\end{center}

\begin{cuadroadvertencia}
\textbf{Cómo fabricar nuevas metas.} Usa enteros entre $-10$ y $+10$ y comprueba que la meta sea alcanzable en pocos turnos con los valores del dado. Evita metas que solo se alcancen con un valor exacto: el juego debe permitir varios caminos.
\end{cuadroadvertencia}

\subsubsection{Patrón C — La maniobra de creación de vacío}

Para extraer $k$ tapas de un color del que solo hay $m<k$:

\begin{enumerate}[leftmargin=*]
  \item Añade $k$ tapas del color contrario y $k$ tapas del color pedido (son $k$ pares nulos).
  \item Extrae las $k$ tapas del color pedido.
  \item Aniquila los pares que queden.
\end{enumerate}

\textbf{Resultado neto:} la carga aumenta en $k$ si el color pedido era rojo (negativo) y disminuye en $k$ si era azul (positivo).

\subsection{Configuración inicial}

\begin{enumerate}[leftmargin=*]
  \item Cada jugador recibe un vaso vacío.
  \item Cada jugador toma \textbf{3 tapas al azar} de una bolsa opaca, las echa en su vaso, aniquila pares y anuncia su carga inicial.
  \item Las tapas restantes forman el \textbf{Banco de Materia y Antimateria}.
  \item Se barajan las 20 Cartas de Meta y se revela la primera.
  \item Se dejan los dos dados y una hoja de borrador por jugador.
\end{enumerate}

% ============================================================
\clearpage
\section{Juego 13 — Vectores de Invasión: la matriz de los signos}
% ============================================================

\subsection{Ficha rápida}

\begin{center}
\small
\begin{tabular}{@{}p{3cm} p{12.2cm}@{}}
\toprule
\textbf{Jugadores} & 2 jugadores o 2 equipos de 2 integrantes\\
\textbf{Duración} & 30--45 minutos\\
\textbf{Objetivo} & Conquistar casillas de una matriz $6\times 6$ multiplicando o dividiendo exactamente los dos números señalados por dos clips.\\
\textbf{Subtemas} & Multiplicación y división exacta en $\Z$, regla de los signos, propiedad distributiva, factor común, elemento absorbente ($0$).\\
\textbf{Qué se evalúa} & Que el estudiante anuncie la operación, justifique el signo y verifique que el cociente sea exacto.\\
\textbf{Nivel} & Profundización.\\
\bottomrule
\end{tabular}
\end{center}

\subsection{La idea en una frase}

\begin{cuadroextra}
  «Dos cursores sobre dos regletas de enteros deciden tu jugada: mueve solo uno, multiplica o divide, y conquista la casilla del resultado. Cuatro en línea o un bloque $2\times 2$ y la matriz es tuya.»
\end{cuadroextra}

\subsection{Cómo se juega en cinco pasos}

\begin{enumerate}[leftmargin=*]
  \item El tablero $6\times 6$ está en el centro; las regletas $A$ y $B$, abajo, con un clip cada una.
  \item El Jugador 1 coloca \textbf{los dos clips} donde quiera, hace la operación y coloca su primera ficha.
  \item En los turnos siguientes, cada jugador mueve \textbf{solo uno} de los dos clips.
  \item Anuncia la operación (producto o cociente exacto), justifica el signo y coloca su ficha en la casilla del resultado.
  \item Gana quien logre \textbf{4 fichas en línea} (horizontal, vertical o diagonal) o un \textbf{bloque $2\times 2$}.
\end{enumerate}

\subsection{Reglas}

\textbf{Idea clave:} los dos clips señalan dos números; en cada turno se mueve uno y se conquista la casilla del resultado de una operación exacta.

\begin{enumerate}[leftmargin=*]
  \item \textbf{Dos regletas y dos clips.} La regleta A y la regleta B tienen 9 valores cada una. Un clip en cada regleta señala el número activo.
  \item \textbf{Primer turno.} El Jugador 1 coloca los dos clips donde quiera y hace su jugada. En los turnos siguientes, cada jugador mueve \textbf{solo uno} de los clips.
  \item \textbf{Movimiento del clip.} El clip puede moverse a cualquier casilla de su propia regleta, siempre que sea distinta de la actual.
  \item \textbf{Operaciones posibles.} Con los dos números señalados se puede hacer $A\cdot B$, $A\div B$ o $B\div A$, siempre que el cociente sea exacto y el divisor sea distinto de $0$.
  \item \textbf{Signo obligatorio.} Antes de colocar la ficha se dice el resultado y se justifica el signo: «$(-3)\cdot(-4)=+12$; menos por menos da más».
  \item \textbf{Casilla ocupada o inexistente.} Si el resultado no está en el tablero o la casilla está ocupada, se prueba otra operación u otro orden. Si ninguna opción sirve, se devuelve \textbf{solo el clip movido} a su posición anterior y se pierde el turno.
  \item \textbf{Cero.} Multiplicar por $0$ permite conquistar una de las dos casillas $0$. \textbf{Dividir entre $0$ está prohibido.}
  \item \textbf{Fichas.} Cada jugador tiene 18 fichas. Si se le acaban, no puede colocar más; solo puede ganar por desempate si el tablero se llena.
  \item \textbf{Victoria inmediata.} Gana quien alinee 4 fichas propias (horizontal, vertical o diagonal) o forme un bloque $2\times 2$. La partida termina en ese momento.
  \item \textbf{Matriz llena.} Si se llenan las 36 casillas sin ganador, cada jugador suma los valores absolutos de las esquinas que tenga ocupadas con sus propias fichas; gana la suma mayor. Si persiste el empate, gana quien tenga más casillas.
  \item \textbf{Carta de Hackeo.} Cada jugador tiene 2 cartas. Una carta se usa una sola vez, como máximo dos por partida y nunca dos en el mismo turno; sustituye el movimiento normal y se resuelve según el apartado «Hackeo Distributivo».
\end{enumerate}

\subsection{Turno paso a paso}

\begin{enumerate}[leftmargin=*]
  \item \textbf{Observa} la posición de los dos clips.
  \item \textbf{Elige} cuál mover y a qué valor de su regleta.
  \item \textbf{Prueba} las operaciones posibles: $A\cdot B$, $A\div B$, $B\div A$.
  \item \textbf{Verifica} que el cociente sea exacto y que la casilla esté libre.
  \item \textbf{Anuncia} la operación y el signo.
  \item \textbf{Coloca} la ficha. Si no hay jugada legal, devuelve el clip y pierde el turno.
  \item \textbf{Comprueba} si lograste 4 en línea o un bloque $2\times 2$.
\end{enumerate}

\subsection{Hackeo Distributivo (carta especial)}

Cada jugador recibe \textbf{2 cartas}. Cada carta se usa \textbf{una sola vez} (máximo 2 usos por partida y nunca dos en el mismo turno) y sustituye el movimiento normal del turno.

\begin{cuadropasos}
\begin{enumerate}[leftmargin=*]
  \item Elige uno de los dos números señalados; llámalo $a$. El otro es $d$.
  \item Descompón $a=b+c$, con $b$ y $c$ enteros \textbf{distintos de cero} y $b\ne a$.
  \item Escribe y resuelve en la hoja de auditoría:
  \[
    d\cdot a=d\cdot(b+c)=d\cdot b+d\cdot c.
  \]
  \item Si el resultado final es una casilla libre, la conquistas. Si no, pierdes la carta y el turno.
\end{enumerate}
\end{cuadropasos}

\textbf{Ejemplo.} Clips en $A=-4$ y $B=+3$. El jugador elige $a=-4$ y lo descompone como $-4=-5+1$:
\[
3\cdot(-4)=3\cdot(-5)+3\cdot(1)=-15+3=-12.
\]
Si la casilla $-12$ está libre, la conquista.

\subsection{Final y puntuación}

\begin{center}
\small
\begin{tabular}{@{}p{7.2cm} p{6.8cm}@{}}
\toprule
\textbf{Situación} & \textbf{Resultado}\\
\midrule
Alinea 4 fichas (horizontal, vertical o diagonal) & Gana de inmediato.\\
Forma un bloque $2\times 2$ & Gana de inmediato.\\
Se llenan las 36 casillas sin ganador & Suma los valores absolutos de sus cuatro esquinas; gana la suma mayor.\\
Persiste el empate & Gana quien tenga más casillas ocupadas.\\
\bottomrule
\end{tabular}
\end{center}

Ganar no da puntos: la partida termina en el momento en que se logra la línea o el bloque.

\subsection{Ejemplo resuelto}

Clips en $A=-12$ y $B=+3$.

\begin{itemize}[leftmargin=*]
  \item Producto: $(-12)\cdot 3=-36$. Si la casilla $-36$ está libre, puede conquistarla.
  \item Cociente: $(-12)\div 3=-4$. También es exacto.
  \item Cociente inverso: $3\div(-12)$ no es exacto; no vale.
\end{itemize}

El jugador mueve el clip $B$ de $+3$ a $-2$ (solo mueve un clip). Ahora:
\[
(-12)\div(-2)=+6.
\]
Anuncia: «$-12$ entre $-2$ es $+6$, porque menos entre menos da más». Si $+6$ está libre, coloca la ficha.

\subsection{Materiales}

\begin{center}
\small
\begin{tabular}{@{}p{5.0cm} c p{7.4cm}@{}}
\toprule
\textbf{Material} & \textbf{Cantidad} & \textbf{Uso}\\
\midrule
Cartulina gruesa o cartón reciclado & $30\times 30$ cm & Tablero y regletas\\
Clips metálicos & 2 & Cursores\\
Fichas verdes & 18 & Jugador 1\\
Fichas naranjas & 18 & Jugador 2\\
Tarjetas de $6\times 9$ cm & 4 & Cartas de Hackeo\\
Marcadores de punta fina y gruesa & 3 & Rotular\\
Regla, escuadra, tijera, lápiz & 1 cada uno & Trazado\\
\bottomrule
\end{tabular}
\end{center}

\subsection{Construcción paso a paso}

\begin{enumerate}[leftmargin=*]
  \item \textbf{Matriz.} Dibuja un cuadrado de $24\times 24$ cm y divídelo en una cuadrícula $6\times 6$ (casillas de $4\times 4$ cm).
  \item \textbf{Regletas.} Debajo del tablero dibuja dos tiras de $2{,}5$ cm de alto, divididas en 9 casillas cada una (o recórtalas aparte para deslizar los clips).
  \item \textbf{Regleta $A$.} De izquierda a derecha: $-12,\ -6,\ -4,\ -3,\ -2,\ 0,\ +2,\ +4,\ +6$.
  \item \textbf{Regleta $B$.} De izquierda a derecha: $-6,\ -4,\ -3,\ -2,\ -1,\ +1,\ +2,\ +3,\ +5$.
  \item \textbf{Matriz.} Copia la cuadrícula del Patrón B, celda por celda.
  \item \textbf{Fichas.} Recorta 18 círculos verdes y 18 naranjas de $2{,}5$ cm de diámetro.
\end{enumerate}

\subsection{Patrones explícitos (fabricación)}

\subsubsection{Patrón A — Las dos regletas}

\begin{center}
\small
\begin{tabular}{@{}c l@{}}
\toprule
\textbf{Regleta} & \textbf{Valores (9 casillas)}\\
\midrule
$A$ & $-12,\ -6,\ -4,\ -3,\ -2,\ 0,\ +2,\ +4,\ +6$\\
$B$ & $-6,\ -4,\ -3,\ -2,\ -1,\ +1,\ +2,\ +3,\ +5$\\
\bottomrule
\end{tabular}
\end{center}

\subsubsection{Patrón B — La matriz $6\times 6$ completa}

\begin{center}
\small
\begin{tabular}{@{}c|r|r|r|r|r|r@{}}
\toprule
 & \textbf{C1} & \textbf{C2} & \textbf{C3} & \textbf{C4} & \textbf{C5} & \textbf{C6}\\
\midrule
\textbf{F1} & $-36$ & $+10$ & $-15$ & $+24$ & $-4$ & $+8$\\
\midrule
\textbf{F2} & $+18$ & $-8$ & $0$ & $-30$ & $+12$ & $-1$\\
\midrule
\textbf{F3} & $-20$ & $+6$ & $+72$ & $+2$ & $-16$ & $+36$\\
\midrule
\textbf{F4} & $+9$ & $-12$ & $+20$ & $-6$ & $+48$ & $+4$\\
\midrule
\textbf{F5} & $-3$ & $+16$ & $-10$ & $+30$ & $-2$ & $-24$\\
\midrule
\textbf{F6} & $0$ & $-18$ & $+3$ & $-9$ & $+1$ & $-60$\\
\bottomrule
\end{tabular}
\end{center}

\textbf{Propiedades de esta matriz:}
\begin{itemize}[leftmargin=*]
  \item Contiene los \textbf{35 resultados distintos} que se pueden formar con las dos regletas.
  \item El $0$ aparece \textbf{dos veces} (F2-C3 y F6-C1): hay dos formas de bloquear con el elemento absorbente.
  \item Las cuatro esquinas son $-36$, $+8$, $0$ y $-60$; se usan en el desempate.
\end{itemize}

\subsubsection{Patrón C — Clave de operaciones (una por casilla)}

\begin{longtable}{@{}r l r l@{}}
\toprule
\textbf{Valor} & \textbf{Operación válida} & \textbf{Valor} & \textbf{Operación válida}\\
\midrule
\endhead
\bottomrule
\endfoot
$-60$ & $-12\cdot 5$ & $+1$ & $-6\div(-6)$\\
$-36$ & $-12\cdot 3$ & $+2$ & $-12\div(-6)$\\
$-30$ & $-6\cdot 5$ & $+3$ & $6\div 2$\\
$-24$ & $-12\cdot 2$ & $+4$ & $2\cdot 2$\\
$-20$ & $-4\cdot 5$ & $+6$ & $2\cdot 3$\\
$-18$ & $-6\cdot 3$ & $+8$ & $4\cdot 2$\\
$-16$ & $4\cdot(-4)$ & $+9$ & $-3\cdot(-3)$\\
$-15$ & $-3\cdot 5$ & $+10$ & $2\cdot 5$\\
$-12$ & $-4\cdot 3$ & $+12$ & $4\cdot 3$\\
$-10$ & $-2\cdot 5$ & $+16$ & $-4\cdot(-4)$\\
$-9$ & $-3\cdot 3$ & $+18$ & $6\cdot 3$\\
$-8$ & $-4\cdot 2$ & $+20$ & $4\cdot 5$\\
$-6$ & $-3\cdot 2$ & $+24$ & $-6\cdot(-4)$\\
$-4$ & $-2\cdot 2$ & $+30$ & $6\cdot 5$\\
$-3$ & $-6\div 2$ & $+36$ & $-6\cdot(-6)$\\
$-2$ & $-6\div 3$ & $+48$ & $-12\cdot(-4)$\\
$-1$ & $-6\div 6$ & $+72$ & $-12\cdot(-6)$\\
$0$ & $0\cdot 5$ & & \\
\end{longtable}

\begin{cuadroadvertencia}
\textbf{Cómo fabricar otra matriz.} Enumera todos los productos y cocientes exactos de las dos regletas; elimina los repetidos; elige 34 valores y añade dos veces el $0$; reparte los 36 valores evitando que los extremos queden juntos. Cualquier celda que no pertenezca al conjunto de resultados hace el tablero injugable.
\end{cuadroadvertencia}

\subsubsection{Patrón D — Carta de Hackeo Distributivo}

\begin{center}
\small
\begin{tabular}{@{}l p{10cm}@{}}
\toprule
\textbf{Elemento} & \textbf{Regla}\\
\midrule
Cuándo se usa & Una vez por carta (máximo 2 por partida), en lugar del movimiento normal.\\
Qué se descompone & Uno de los dos números señalados: $a=b+c$.\\
Restricciones & $b\ne 0$, $c\ne 0$, $b\ne a$.\\
Qué se escribe & $d\cdot a=d\cdot(b+c)=d\cdot b+d\cdot c$.\\
Qué se conquista & La casilla del resultado final, si está libre.\\
Si falla & Se pierde la carta y el turno.\\
\bottomrule
\end{tabular}
\end{center}

\subsection{Configuración inicial}

\begin{enumerate}[leftmargin=*]
  \item Coloca el tablero en el centro.
  \item Sitúa las regletas $A$ y $B$ al pie del tablero, con un clip fuera de las casillas.
  \item Entrega 18 fichas de un color a cada jugador y 2 cartas de Hackeo a cada uno.
  \item Deja una hoja de borrador por jugador para auditorías y hackeos.
\end{enumerate}

% ============================================================
\clearpage
\section{Juego 14 — Cortocircuito $\Z$: jerarquía de voltaje}
% ============================================================

\subsection{Ficha rápida}

\begin{center}
\small
\begin{tabular}{@{}p{3cm} p{12.2cm}@{}}
\toprule
\textbf{Jugadores} & 2 a 4\\
\textbf{Duración} & 25--35 minutos\\
\textbf{Objetivo} & Resolver operaciones combinadas \textbf{sin signos de agrupación} respetando la jerarquía y auditar las jugadas rivales.\\
\textbf{Subtemas} & Adición, sustracción, multiplicación y división exacta en $\Z$; prioridad de operaciones; izquierda a derecha entre operaciones del mismo nivel.\\
\textbf{Qué se evalúa} & Que el estudiante resuelva primero productos y cocientes y justifique el orden.\\
\textbf{Nivel} & Profundización.\\
\bottomrule
\end{tabular}
\end{center}

\subsection{La idea en una frase}

\begin{cuadroextra}
  «Cada loseta es un tramo de red eléctrica: a la izquierda entra un voltaje y a la derecha sale un circuito. Si calculas mal la jerarquía, la red hace cortocircuito y los auditores te lo cobran.»
\end{cuadroextra}

\subsection{Cómo se juega en cinco pasos}

\begin{enumerate}[leftmargin=*]
  \item Se reparten las losetas y se destapa una para iniciar la red.
  \item En tu turno, mira los dos extremos abiertos de la cadena.
  \item Elige una loseta cuya \textbf{Entrada} coincida con el \textbf{valor} de una expresión abierta (o cuya expresión valga el número abierto).
  \item Coloca la loseta y da \textbf{10 segundos} para que los rivales auditen.
  \item Si nadie grita «¡Cortocircuito!», pasa el turno. Gana quien coloque su última loseta.
\end{enumerate}

\subsection{Reglas}

\textbf{Idea clave:} las losetas forman una cadena; cada jugador calcula en silencio y los rivales auditan la conexión.

\begin{enumerate}[leftmargin=*]
  \item \textbf{Jerarquía inquebrantable.} Primero multiplicación y división, de izquierda a derecha; después adición y sustracción, de izquierda a derecha.
  \item \textbf{Orientación fija.} Las losetas no se giran: la \textbf{Entrada} (número) queda a la izquierda y el \textbf{Circuito de Salida} (expresión), a la derecha.
  \item \textbf{Conexión por la derecha.} La Entrada de la nueva loseta debe ser igual al valor de la expresión abierta del extremo derecho.
  \item \textbf{Conexión por la izquierda.} La expresión de la nueva loseta debe valer el número abierto del extremo izquierdo; la loseta se coloca a la izquierda.
  \item \textbf{Conexión silenciosa.} El jugador \textbf{no dice} el resultado en voz alta: lo calcula en su borrador y coloca la loseta. La auditoría es de los rivales.
  \item \textbf{Auditoría «¡Cortocircuito!».} Cualquier rival puede gritarlo dentro de los 10 segundos. Se abre la \textbf{Hoja de Multímetro} y ambos resuelven paso a paso:
  \begin{itemize}
    \item si la conexión era \textbf{incorrecta}, el jugador retira su loseta, roba \textbf{2 losetas de penalización} (si quedan en el pozo) y el auditor gana \textbf{1 ficha de Fusible};
    \item si la conexión era \textbf{correcta} (falsa alarma), el auditor roba \textbf{2 losetas de penalización}.
  \end{itemize}
  \item \textbf{Sin jugada.} Si no tienes loseta que conecte, robas de a una del \textbf{Generador Reserva} hasta poder jugar o hasta que se agote. Si el pozo está vacío y no puedes jugar, pasas sin perder ficha.
  \item \textbf{Sin correcciones tardías.} Una loseta colocada no se retira ni se cambia, salvo que prospere una auditoría.
  \item \textbf{Victoria.} La partida termina cuando un jugador coloca su última loseta («¡Red Estable!») o cuando nadie puede conectar. Si se bloquea, cada jugador cuenta sus losetas restantes y resta 1 por cada ficha de Fusible; gana quien tenga el menor total.
  \item \textbf{Empate.} Si hay empate en losetas restantes, gana quien tenga más fichas de Fusible.
\end{enumerate}

\subsection{Turno paso a paso}

\begin{enumerate}[leftmargin=*]
  \item \textbf{Mira} los dos extremos abiertos: uno es un número y el otro, una expresión.
  \item \textbf{Evalúa} mentalmente la expresión abierta.
  \item \textbf{Busca} en tu mano una loseta compatible.
  \item \textbf{Colócala} sin decir el resultado.
  \item \textbf{Espera} 10 segundos de auditoría.
  \item \textbf{Si hay reclamo}, resuelve en la Hoja de Multímetro y aplica la sanción o la falsa alarma.
  \item \textbf{Roba o pasa} según corresponda.
  \item \textbf{Comprueba} si colocaste tu última loseta o si la cadena quedó bloqueada.
\end{enumerate}

\subsection{Final y puntuación}

\begin{center}
\small
\begin{tabular}{@{}p{7.2cm} p{6.8cm}@{}}
\toprule
\textbf{Situación} & \textbf{Resultado}\\
\midrule
Un jugador coloca su última loseta & Gana («¡Red Estable!»).\\
La cadena se bloquea y nadie puede conectar & Gana quien tenga menos losetas restantes, restando 1 por cada ficha de Fusible.\\
Empate en losetas restantes & Gana quien tenga más fichas de Fusible.\\
Sin losetas en la mano y sin jugada posible & Pasa el turno; no pierde por ello.\\
\bottomrule
\end{tabular}
\end{center}

\subsection{Ejemplo resuelto (con trampa incluida)}

En la mesa está abierto el número $-7$. Un jugador coloca la loseta:
\[
\text{Entrada }[6]\quad|\quad\text{Salida }[\,9+4\cdot(-1)\,].
\]
\begin{itemize}[leftmargin=*]
  \item \textbf{Resolución correcta:} primero $4\cdot(-1)=-4$; luego $9+(-4)=+5$. La loseta vale $5$, no $-7$: \textbf{la conexión es incorrecta}.
  \item \textbf{Trampa:} quien resuelve de izquierda a derecha sin jerarquía hace $(9+4)\cdot(-1)=-13$. Ese $-13$ sí es un número del juego (es la Entrada de otra loseta), así que el error «encaja» y solo se descubre auditando.
\end{itemize}

\subsection{Materiales}

\begin{center}
\small
\begin{tabular}{@{}p{5.0cm} c p{7.4cm}@{}}
\toprule
\textbf{Material} & \textbf{Cantidad} & \textbf{Uso}\\
\midrule
Cartón reciclado o cartulina doble faz & 2 cajas o 2 pliegos & 28 losetas\\
Hojas blancas recicladas & 2 & Forrar el cartón\\
Marcadores & negro, rojo, verde & Rotular\\
Hojas de reciclaje y lápices & 1 por jugador & Hoja de Multímetro\\
Regla, tijera, pega & 1 cada uno & Armado\\
\bottomrule
\end{tabular}
\end{center}

\subsection{Construcción paso a paso}

\begin{enumerate}[leftmargin=*]
  \item \textbf{Forrado.} Pega las hojas blancas sobre el cartón y deja secar.
  \item \textbf{Corte.} Traza 28 rectángulos de $4\times 8$ cm y recórtalos.
  \item \textbf{División.} Traza una línea vertical a los $4$ cm: lado izquierdo, ícono de enchufe (Entrada); lado derecho, ícono de rayo (Salida).
  \item \textbf{Rotulado.} Copia la tabla del Patrón B: cada loseta lleva su Entrada y su expresión.
  \item \textbf{Control de calidad.} Verifica que el resultado de cada loseta sea la Entrada de la siguiente (incluida la última, que cierra con la primera).
\end{enumerate}

\subsection{Patrones explícitos (fabricación)}

\subsubsection{Patrón A — Regla del bucle cerrado}

El mazo es \textbf{un ciclo de 28 losetas}:
\[
\text{Salida de la loseta }k=\text{Entrada de la loseta }k+1,\qquad
\text{Salida de la 28}=\text{Entrada de la 1}.
\]
Para fabricar otro mazo: elige 28 resultados distintos entre $-20$ y $+20$, asigna una expresión a cada uno y ordena las losetas de modo que cada Salida sea la Entrada de la siguiente.

\subsubsection{Patrón B — Las 28 losetas (mazo completo)}

{\small
\begin{longtable}{@{}r c l c c@{}}
\toprule
\textbf{\#} & \textbf{Entrada} & \textbf{Circuito de Salida} & \textbf{Resultado} & \textbf{Trampa}\\
\midrule
\endhead
\bottomrule
\endfoot
1 & $+1$ & $2-1+(-1)\cdot(-1)$ & $+2$ & $0$\\
2 & $+2$ & $-10-(-7)\cdot 4$ & $+18$ & $-12$\\
3 & $+18$ & $-3-3+3\cdot(-4)$ & $-18$ & $+12$\\
4 & $-18$ & $-6+12\div 3$ & $-2$ & $+2$\\
5 & $-2$ & $3+3\cdot(-3)$ & $-6$ & $-18$\\
6 & $-6$ & $-5-(-2)+4\cdot 3$ & $+9$ & $+3$\\
7 & $+9$ & $11-1\cdot(-4)$ & $+15$ & $-40$\\
8 & $+15$ & $5-(-3)\cdot 2$ & $+11$ & $+16$\\
9 & $+11$ & $3-(-1)\cdot 3$ & $+6$ & $+12$\\
10 & $+6$ & $9+4\cdot(-1)$ & $+5$ & $-13$\\
11 & $+5$ & $0-(-4)+(-4)\cdot(-1)$ & $+8$ & $0$\\
12 & $+8$ & $-3-(-4)\cdot 5$ & $+17$ & $+5$\\
13 & $+17$ & $-11+9\cdot 2$ & $+7$ & $-4$\\
14 & $+7$ & $7-(-2)+4\cdot(-3)$ & $-3$ & $-39$\\
15 & $-3$ & $-4-8\div 2$ & $-8$ & $-6$\\
16 & $-8$ & $-12-(-8)\div 4$ & $-10$ & $-1$\\
17 & $-10$ & $-1+(-6)\cdot 3$ & $-19$ & $-21$\\
18 & $-19$ & $6+9\div(-3)$ & $+3$ & $-5$\\
19 & $+3$ & $5+4\cdot(-5)$ & $-15$ & $-45$\\
20 & $-15$ & $2-(-6)\cdot(-1)$ & $-4$ & $-8$\\
21 & $-4$ & $11+(-6)\cdot 4$ & $-13$ & $+20$\\
22 & $-13$ & $8-(-1)\cdot 4$ & $+12$ & $+36$\\
23 & $+12$ & $-1-2+(-2)\cdot 3$ & $-9$ & $-15$\\
24 & $-9$ & $-8-(-6)\div 2$ & $-5$ & $-1$\\
25 & $-5$ & $-3+6\div 3$ & $-1$ & $+1$\\
26 & $-1$ & $1+(-3)\cdot 5$ & $-14$ & $-10$\\
27 & $-14$ & $-3-7+(-1)\cdot 2$ & $-12$ & $-22$\\
28 & $-12$ & $2+(-2)\div 2$ & $+1$ & $0$\\
\end{longtable}
}

\textbf{Verificación rápida:} el Resultado de la loseta 28 es $+1$, que es la Entrada de la loseta 1. El bucle cierra.

\subsubsection{Patrón C — Tipos de trampa de jerarquía}

\begin{center}
\small
\begin{tabular}{@{}c l l l@{}}
\toprule
\textbf{Código} & \textbf{Expresión} & \textbf{Error típico} & \textbf{Resultado correcto}\\
\midrule
T1 & $a+b\cdot c$ & $(a+b)\cdot c$ & $a+bc$\\
T2 & $a-b\cdot c$ & $(a-b)\cdot c$ & $a-bc$\\
T3 & $a+b\div c$ & $(a+b)\div c$ & $a+b/c$\\
T4 & $a-b\div c$ & $(a-b)\div c$ & $a-b/c$\\
T5 & $a-b+c\cdot d$ & $(a-b+c)\cdot d$ & $a-b+cd$\\
\bottomrule
\end{tabular}
\end{center}

\begin{cuadroadvertencia}
\textbf{Cómo fabricar nuevas losetas.} Elige un tipo de trampa, números pequeños (entre $-12$ y $+12$), comprueba que todas las divisiones sean exactas y que la trampa no coincida con el resultado correcto. Después encadena la Salida con la Entrada de otra loseta.
\end{cuadroadvertencia}

\subsubsection{Patrón D — Mecánica de auditoría}

\begin{center}
\small
\begin{tabular}{@{}l p{9.6cm}@{}}
\toprule
\textbf{Situación} & \textbf{Consecuencia}\\
\midrule
Nadie reclama en 10 s & La jugada vale y pasa el turno.\\
Reclamo y la conexión era incorrecta & El jugador retira la loseta, roba 2 de penalización; el auditor gana 1 Fusible.\\
Reclamo y la conexión era correcta & El auditor roba 2 de penalización.\\
Se acaba el pozo y no hay jugada & El jugador pasa.\\
\bottomrule
\end{tabular}
\end{center}

\subsection{Configuración inicial}

\begin{enumerate}[leftmargin=*]
  \item Coloca las 28 losetas boca abajo y revuélvelas.
  \item Reparte \textbf{5 losetas} a cada jugador (7 si son solo 2 jugadores).
  \item Las sobrantes forman el \textbf{Generador Reserva}.
  \item Destapa una loseta en el centro para inaugurar la red.
  \item Entrega media hoja de borrador («Multímetro») y un lápiz a cada jugador.
\end{enumerate}

% ============================================================
\clearpage
\section{Juego 15 — Código Bóveda: desactivación concéntrica}
% ============================================================

\subsection{Ficha rápida}

\begin{center}
\small
\begin{tabular}{@{}p{3cm} p{12.2cm}@{}}
\toprule
\textbf{Jugadores} & Cooperativo (2--4) o duelo por equipos (2 vs. 2)\\
\textbf{Duración} & 30--45 minutos\\
\textbf{Objetivo} & Resolver operaciones combinadas con signos de agrupación anidados, de adentro hacia afuera, girando tres anillos concéntricos.\\
\textbf{Subtemas} & Paréntesis $(~)$, corchetes $[~]$, llaves $\{~\}$, valor absoluto $\abs{~}$; supresión de signos de agrupación precedidos por $+$, por $-$ o por un factor.\\
\textbf{Qué se evalúa} & Que el estudiante respete el orden de las capas y el cambio de signo delante de un grupo.\\
\textbf{Nivel} & Profundización.\\
\bottomrule
\end{tabular}
\end{center}

\subsection{La idea en una frase}

\begin{cuadroextra}
  «La bóveda tiene tres anillos: paréntesis, corchetes y llaves. Solo se abre resolviendo desde el núcleo hacia afuera y girando cada anillo hasta el número exacto.»
\end{cuadroextra}

\subsection{Cómo se juega en cinco pasos}

\begin{enumerate}[leftmargin=*]
  \item Se coloca el disco de tres anillos en el centro, alineado en el $0$.
  \item Se revela una \textbf{Tarjeta de Bóveda} y arranca el cronómetro ($3$ min).
  \item El \textbf{Operador de Paréntesis} resuelve el núcleo y gira el anillo central.
  \item El \textbf{Operador de Corchetes} sustituye ese valor, resuelve y gira el anillo medio.
  \item El \textbf{Operador de Llaves} hace lo mismo con el anillo exterior y el equipo calcula la \textbf{Clave Final}. Si los tres anillos y la clave coinciden, la bóveda se abre.
\end{enumerate}

\subsection{Reglas}

\textbf{Idea clave:} se resuelve desde el núcleo hacia afuera y cada anillo se gira al valor parcial de su capa, no al resultado final.

\begin{enumerate}[leftmargin=*]
  \item \textbf{De adentro hacia afuera.} No se mueve el anillo de corchetes sin haber fijado el de paréntesis; no se mueve el de llaves sin haber fijado el de corchetes.
  \item \textbf{Qué muestra cada anillo.} El anillo central muestra el valor del paréntesis; el medio, el del corchete; el exterior, el de la llave.
  \item \textbf{Cambio de polaridad.} Si un signo de agrupación está precedido por un signo $-$, el operador de esa capa anuncia «¡Inversión de polaridad!» y cambia el signo de \textbf{todos} los términos internos (o toma el opuesto del resultado del grupo).
  \item \textbf{Valor absoluto.} Las barras $\abs{~}$ son una capa más: primero se calcula su contenido y luego se toma el valor no negativo.
  \item \textbf{Roles.} En equipos de 3 o 2, cada integrante es responsable de una capa; el cuarto integrante, si lo hay, es \textbf{verificador} y revisa la Clave Final.
  \item \textbf{Tiempo.} Cada bóveda dura $3$ minutos. Al agotarse, se considera fallo y se suma una alarma.
  \item \textbf{Cascada.} Si el anillo central es incorrecto, los anillos exteriores quedan automáticamente inválidos: el error se arrastra. Solo se puede corregir una capa antes de girar la siguiente.
  \item \textbf{Entrega.} El verificador anuncia la Clave Final y el equipo la escribe en el Bloc de Criptografía antes de voltear la tarjeta.
  \item \textbf{Auditoría.} Al terminar, el equipo voltea la tarjeta y compara los tres valores parciales y la clave final. La tarjeta muestra el «Código Triple de Alineación».
  \item \textbf{Fin.} El modo cooperativo y el competitivo se puntúan como se indica en «Final y puntuación».
\end{enumerate}

\subsection{Turno paso a paso}

\begin{enumerate}[leftmargin=*]
  \item \textbf{Lee} la expresión completa y localiza la capa más interna: $(~)$.
  \item \textbf{Resuelve} el paréntesis respetando la jerarquía interna.
  \item \textbf{Gira} el anillo central hasta ese valor.
  \item \textbf{Sustituye} el valor en el corchete, resuelve y gira el anillo medio.
  \item \textbf{Sustituye} el valor en la llave, resuelve y gira el anillo exterior.
  \item \textbf{Calcula} las operaciones externas (factores y valor absoluto) y obtén la Clave Final.
  \item \textbf{Verifica} con el dorso de la tarjeta.
  \item \textbf{Anota} el resultado y pasa a la bóveda siguiente.
\end{enumerate}

\subsection{Final y puntuación}

\subsubsection{Modo cooperativo}

El equipo debe abrir \textbf{4 bóvedas} (Bronce, Plata, Oro y Diamante) antes de acumular \textbf{3 alarmas}.

\begin{center}
\small
\begin{tabular}{@{}l c p{9cm}@{}}
\toprule
\textbf{Nivel} & \textbf{Dificultad} & \textbf{Contenido}\\
\midrule
Bronce & 1 estrella & Solo adición y sustracción con paréntesis y corchetes.\\
Plata & 1--2 estrellas & Paréntesis, corchetes y un signo negativo delante de un grupo.\\
Oro & 2 estrellas & Multiplicación o división delante de un grupo.\\
Diamante & 3 estrellas & Llaves anidadas y valor absoluto.\\
\bottomrule
\end{tabular}
\end{center}

Cada fallo (código incorrecto o tiempo agotado) suma \textbf{1 alarma}. Con \textbf{3 alarmas}, la bóveda queda sellada y el equipo pierde.

\subsubsection{Modo competitivo (2 vs. 2)}

Se juegan \textbf{5 rondas simultáneas}; cada equipo tiene su propio disco y recibe una tarjeta del mismo nivel.

\begin{center}
\small
\begin{tabular}{@{}l c@{}}
\toprule
\textbf{Logro en la ronda} & \textbf{Puntos}\\
\midrule
Cada anillo correctamente posicionado, con las capas interiores correctas & $1$ (máximo $3$)\\
Clave Final correcta & $2$\\
\textbf{Máximo por ronda} & $5$\\
\bottomrule
\end{tabular}
\end{center}

Gana el equipo con más puntos tras las 5 rondas. Desempate: (a) más Claves Finales exactas; (b) menor tiempo total.

\subsection{Ejemplo resuelto}

Tarjeta:
\[
-2\cdot\{5-[3+(-8+12\div(-3))-(-2)]\}-\abs{-6}
\]

\begin{center}
\small
\begin{tabular}{@{}l l r l@{}}
\toprule
\textbf{Capa} & \textbf{Operación} & \textbf{Valor} & \textbf{Anillo}\\
\midrule
Paréntesis $(~)$ & $-8+12\div(-3)=-8+(-4)$ & $-12$ & Central\\
Corchetes $[~]$ & $3+(-12)-(-2)=3-12+2$ & $-7$ & Medio\\
Llaves $\{~\}$ & $5-(-7)=5+7$ & $+12$ & Exterior\\
Clave Final & $-2\cdot(+12)-\abs{-6}=-24-6$ & $-30$ & ---\\
\bottomrule
\end{tabular}
\end{center}

\textbf{Código Triple de Alineación:} $(-12)\rightarrow[-7]\rightarrow\{+12\}$; \textbf{Clave Final: $-30$}.

\subsection{Materiales}

\begin{center}
\small
\begin{tabular}{@{}p{5.0cm} c p{7.4cm}@{}}
\toprule
\textbf{Material} & \textbf{Cantidad} & \textbf{Uso}\\
\midrule
Cartulinas de colores & 3 (6 si es competitivo) & Un disco por capa y por equipo\\
Chinche de dos patas o botón + clip + pabilo & 1 & Eje giratorio\\
Cartulina blanca & 1 & 12 Tarjetas de Bóveda\\
Marcadores de punta fina y gruesa & 3 & Rotular\\
Reloj de arena o cronómetro & 1 & 3 minutos por bóveda\\
Compás, regla, lápiz, tijera & 1 cada uno & Trazado\\
\bottomrule
\end{tabular}
\end{center}

\subsection{Construcción paso a paso}

\begin{enumerate}[leftmargin=*]
  \item \textbf{Discos.} Traza y recorta tres círculos: exterior (llaves) de $10$ cm de radio, intermedio (corchetes) de $7{,}5$ cm y central (paréntesis) de $5$ cm.
  \item \textbf{Sectores.} Divide el borde de cada disco en \textbf{16 sectores} iguales (cuartos, octavos y dieciseisavos).
  \item \textbf{Numeración.} En el borde de los tres discos escribe los mismos 16 números del Patrón A. En el centro del disco pequeño escribe $(~)$, en la franja visible del mediano $[~]$ y en la del grande $\{~\}$.
  \item \textbf{Eje.} Perfora el centro de los tres discos, apílalos de mayor a menor y atraviesa el chinche (o pasa el pabilo con dos botones) para que giren de forma independiente.
  \item \textbf{Tarjetas.} Recorta 12 tarjetas de $8\times 12$ cm. En la cara frontal va la expresión y el nivel; en el dorso, el código paso a paso.
\end{enumerate}

\subsection{Patrones explícitos (fabricación)}

\subsubsection{Patrón A — Los 16 números de los anillos}

Todos los resultados parciales de las tarjetas deben pertenecer a este conjunto, porque son los únicos números impresos en los discos:
\[
\{-20,-15,-12,-10,-8,-7,-5,-3,\ 0,\ +2,+4,+6,+8,+12,+15,+18\}.
\]

\subsubsection{Patrón B — Las 12 Tarjetas de Bóveda}

{\footnotesize
\begin{longtable}{@{}r c >{\raggedright\arraybackslash}p{6.1cm} c c c c@{}}
\toprule
\textbf{\#} & \textbf{Nivel} & \textbf{Expresión} & $\boldsymbol{(~)}$ & $\boldsymbol{[~]}$ & $\boldsymbol{\{~\}}$ & \textbf{Clave}\\
\midrule
\endhead
\bottomrule
\endfoot
V1 & 1 & $-3+\{12-[5+(2-7)]\}$ & $-5$ & $0$ & $+12$ & $+9$\\
V2 & 1 & $4-\{-10-[2+(2-7)]\}$ & $-5$ & $-3$ & $-7$ & $+11$\\
V3 & 1 & $\{15-[-3-(5-8)]\}+6$ & $-3$ & $0$ & $+15$ & $+21$\\
V4 & 1 & $-8-\{-10-[5+(2-10)]\}$ & $-8$ & $-3$ & $-7$ & $-1$\\
V5 & 2 & $2\cdot\{2+[4\cdot(-2)-(9\div(-3))]\}$ & $-3$ & $-5$ & $-3$ & $-6$\\
V6 & 2 & $-3\cdot\{-3+[(12\div(-4))\cdot 2-1]\}$ & $-3$ & $-7$ & $-10$ & $+30$\\
V7 & 2 & $\{8-[(-15\div 5)\cdot(-2)]\}\div 2$ & $-3$ & $+6$ & $+2$ & $+1$\\
V8 & 2 & $2\cdot\{-6-[(-24\div 3)-4]\}$ & $-8$ & $-12$ & $+6$ & $+12$\\
V9 & 3 & $-2\cdot\{5-[3+(-8+12\div(-3))-(-2)]\}-\abs{-6}$ & $-12$ & $-7$ & $+12$ & $-30$\\
V10 & 3 & $\abs{-3}\cdot\{-6+[-6-(2-14)]\}$ & $-12$ & $+6$ & $0$ & $0$\\
V11 & 3 & $-\abs{-5}+\{15-[(-4)\cdot 3-(4-16)]\}$ & $-12$ & $0$ & $+15$ & $+10$\\
V12 & 3 & $\{[((-6)\cdot 2)\div(-2)]\div 3\}\cdot\abs{-7}$ & $-12$ & $+6$ & $+2$ & $+14$\\
\end{longtable}
}

\subsubsection{Patrón C — Cómo diseñar una tarjeta nueva}

\begin{enumerate}[leftmargin=*]
  \item Elige un valor final cualquiera.
  \item Construye de afuera hacia adentro: decide el valor de las llaves ($B$), luego el de los corchetes ($b$) y luego el de los paréntesis ($p$), \textbf{los tres dentro del conjunto de 16 números}.
  \item Escribe cada capa como una operación que produzca exactamente el valor elegido.
  \item Añade, si es nivel 2 o 3, un factor o un valor absoluto externo.
  \item Verifica el código: $p\rightarrow b\rightarrow B\rightarrow$ Clave Final.
\end{enumerate}

\subsubsection{Patrón D — Cambio de polaridad}

\begin{center}
\small
\begin{tabular}{@{}l l l@{}}
\toprule
\textbf{Expresión} & \textbf{Error frecuente} & \textbf{Procedimiento correcto}\\
\midrule
$-\{5-[3-8]\}$ & Cambiar solo el primer término & $-\{5-(-5)\}=-\{10\}=-10$\\
$-[\,-7+2\,]$ & Dejar el $2$ con el mismo signo & $7-2=+5$\\
$-2\cdot\{5-[-7]\}$ & Multiplicar solo el $5$ & $-2\cdot\{12\}=-24$\\
\bottomrule
\end{tabular}
\end{center}

\begin{cuadroextra}
\textbf{Regla mnemotécnica:} «El menos de afuera cambia el signo de todo lo de adentro; el factor de afuera multiplica todo el grupo, no solo el primer término».
\end{cuadroextra}

\subsection{Configuración inicial}

\begin{enumerate}[leftmargin=*]
  \item Coloca el disco con los tres anillos alineados en el $0$ («Bóveda Bloqueada»).
  \item Separa las tarjetas por nivel: Nivel 1 (V1--V4), Nivel 2 (V5--V8) y Nivel 3 (V9--V12).
  \item Entrega una hoja de cálculo («Bloc de Criptografía») por jugador o equipo.
  \item Prepara el cronómetro en $3$ minutos.
  \item En modo competitivo, entrega un disco y un juego de tarjetas a cada equipo.
\end{enumerate}

% ============================================================
\appendix
% ============================================================
\clearpage

\section{Resumen comparativo}

\begin{center}
\footnotesize
\begin{tabular}{@{}c l p{3.0cm} p{3.4cm} p{3.6cm}@{}}
\toprule
\textbf{\#} & \textbf{Juego} & \textbf{Subtemas} & \textbf{Mecánica} & \textbf{Señal de aprendizaje}\\
\midrule
1 & Semáforo de enteros & Orden, opuestos, valor absoluto, operaciones & Verdadero o falso con tarjetas & Elige y justifica la regla.\\
2 & El termómetro humano & Ubicación y orden & Fila humana sobre una cinta & Ubica y nombra un mayor y un menor.\\
3 & El ascensor de dados & Suma y resta de una operación & Dado de cargas y ficha en una tira & Anuncia la operación antes de mover.\\
4 & Pesca de opuestos & Números opuestos & Pesca de cartas entre compañeros & Nombra el opuesto y explica que suma cero.\\
5 & Tres en raya de enteros & Suma y resta de una operación & Cuadrícula $3\times 3$ con resultados & Resuelve y marca el resultado.\\
6 & Guerra de pisos & Orden y comparación & Guerra de cartas con la recta vertical & Justifica la comparación con la posición.\\
7 & Memoria de opuestos & Opuestos y valor absoluto & Memoria de parejas & Explica por qué dos cartas son equivalentes.\\
8 & La serpiente de los enteros & Suma y resta de una operación & Saltos en un tablero de 30 casillas & Decide si avanza o retrocede y por qué.\\
9 & Bingo de productos & Producto, cociente y signos & Bingo $4\times 4$ con dos mazos & Justifica el signo del resultado.\\
10 & La escalera de dos peldaños & Combinadas de dos pasos & Carrera de 8 peldaños & Resuelve primero el producto o cociente.\\
11 & Abismo y Cumbre & Orden, opuestos, valor absoluto & Recta numérica y captura por simetría & Justifica por qué dos números son simétricos.\\
12 & Laboratorio de Antimateria & Adición, sustracción, propiedades & Reactor de cargas y pares nulos & Explica por qué restar un negativo suma.\\
13 & Vectores de Invasión & Producto, cociente, signos, distributiva & Conquista de matriz $6\times 6$ & Anuncia la operación y justifica el signo.\\
14 & Cortocircuito $\Z$ & Combinadas sin agrupación & Dominó con auditoría & Resuelve primero productos y cocientes.\\
15 & Código Bóveda & Combinadas con agrupación & Tres anillos concéntricos & Respeta las capas y el cambio de signo.\\
\bottomrule
\end{tabular}
\end{center}

\section{Todos los patrones en un solo lugar}

\begin{longtable}{@{}c l l@{}}
\toprule
\textbf{Juego} & \textbf{Patrón} & \textbf{Dónde está}\\
\midrule
\endhead
\bottomrule
\endfoot
1 & 24 afirmaciones de verdadero/falso & Sección 1.10, Patrón A\\
1 & Regla de puntuación & Sección 1.10, Patrón B\\
1 & Cómo fabricar nuevas afirmaciones & Sección 1.10, Patrón C\\
2 & 20 tarjetas de enteros & Sección 2.10, Patrón A\\
2 & Disposición de la cinta & Sección 2.10, Patrón B\\
2 & Frases de comparación & Sección 2.10, Patrón C\\
3 & Caras del dado de cargas & Sección 3.10, Patrón A\\
3 & Tira numérica de $-10$ a $+10$ & Sección 3.10, Patrón B\\
3 & Regla de movimiento y rebote & Sección 3.10, Patrón C\\
3 & Variante con moneda y dado común & Sección 3.10, Patrón D\\
4 & Composición del mazo (32 cartas) & Sección 4.10, Patrón A\\
4 & Parejas válidas de opuestos & Sección 4.10, Patrón B\\
4 & Guion de la pregunta & Sección 4.10, Patrón C\\
5 & Tablero $3\times 3$ & Sección 5.10, Patrón A\\
5 & 27 tarjetas de suma y resta & Sección 5.10, Patrón B\\
5 & Clave por resultado & Sección 5.10, Patrón C\\
6 & Composición del mazo: 20 valores $\times$ 2 & Sección 6.10, Patrón A\\
6 & Las tres justificaciones válidas & Sección 6.10, Patrón B\\
6 & Regla de desempate & Sección 6.10, Patrón C\\
7 & 12 pares de opuestos (24 cartas) & Sección 7.10, Patrón A\\
7 & 8 pares de valor absoluto (16 cartas) & Sección 7.10, Patrón B\\
7 & Cuadrícula de colocación $8\times 5$ & Sección 7.10, Patrón C\\
8 & Las 24 cartas de suma y resta & Sección 8.10, Patrón A\\
8 & Regla del salto y del rebote & Sección 8.10, Patrón B\\
9 & Los dos mazos ($-6$ a $+6$) & Sección 9.10, Patrón A\\
9 & Los 4 cartones de $4\times 4$ & Sección 9.10, Patrón B\\
9 & Clave de operaciones por valor & Sección 9.10, Patrón C\\
10 & Escaleras A, B, C y D & Sección 10.10, Patrones A--D\\
10 & Regla de diseño de una escalera & Sección 10.10, Patrón E\\
11 & P1--P7: tipos de expresión de la cara A & Sección 11.10, Patrón A\\
11 & 29 tarjetas de coordenadas (mazo completo) & Sección 11.10, Patrón B\\
11 & 14 pares de opuestos + la tarjeta $-30$ & Sección 11.10, Patrón C\\
11 & 6 tarjetas de sonda con su conjunto solución & Sección 11.10, Patrón D\\
12 & Caras de los dos dados & Sección 12.10, Patrón A\\
12 & 20 Cartas de Meta & Sección 12.10, Patrón B\\
12 & Maniobra de creación de vacío & Sección 12.10, Patrón C\\
13 & Valores de las regletas $A$ y $B$ & Sección 13.11, Patrón A\\
13 & Matriz $6\times 6$ completa & Sección 13.11, Patrón B\\
13 & Clave de operaciones por casilla & Sección 13.11, Patrón C\\
13 & Hackeo Distributivo & Sección 13.11, Patrón D\\
14 & Regla del bucle cerrado & Sección 14.10, Patrón A\\
14 & 28 losetas completas & Sección 14.10, Patrón B\\
14 & Tipos de trampa T1--T5 & Sección 14.10, Patrón C\\
14 & Auditoría y sanciones & Sección 14.10, Patrón D\\
15 & 16 números de los anillos & Sección 15.10, Patrón A\\
15 & 12 Tarjetas de Bóveda & Sección 15.10, Patrón B\\
15 & Diseño de una tarjeta nueva & Sección 15.10, Patrón C\\
15 & Cambio de polaridad & Sección 15.10, Patrón D\\
\end{longtable}

\section{Plantilla de tarjeta en blanco}

\begin{center}
\small
\begin{tabular}{@{}l p{10.5cm}@{}}
\toprule
\textbf{Campo} & \textbf{Qué escribir}\\
\midrule
Juego & 1 a 10.\\
Nivel / dificultad & 1 a 3 estrellas.\\
Cara A & Expresión, condición o reto (sin el resultado).\\
Cara B & Resultado, solución paso a paso o código de auditoría.\\
Patrón usado & P1--P7, T1--T5, etc.\\
Verificación & Firma del docente que comprobó la operación.\\
\bottomrule
\end{tabular}
\end{center}

\section{Rúbrica de observación}

\begin{center}
\small
\begin{tabular}{@{}p{3.4cm} p{3.6cm} p{3.6cm} p{3.6cm}@{}}
\toprule
\textbf{Indicador} & \textbf{1 (En inicio)} & \textbf{2 (En proceso)} & \textbf{3 (Logrado)}\\
\midrule
Resolución de la operación & Se equivoca y no detecta el error. & Llega al resultado con ayuda. & Resuelve solo y verifica.\\
Justificación & No explica el procedimiento. & Explica parcialmente. & Explica el orden y la regla aplicada.\\
Uso del material & Manipula sin relación con la operación. & Usa el material con apoyo. & Traduce la operación al material con autonomía.\\
Auditoría & No detecta errores ajenos. & Detecta el error pero no lo explica. & Detecta, explica y corrige.\\
Trabajo en equipo & Interrumpe o no participa. & Participa si se le pide. & Cumple su rol y ayuda a los demás.\\
\bottomrule
\end{tabular}
\end{center}

\end{document}
